\documentclass[11pt,a4paper,oneside]{report}

% Only packages that TeX Live (CI) and MiKTeX both ship. tests/test_manual.py checks the list.
\usepackage[T1]{fontenc}
\usepackage[utf8]{inputenc}
\usepackage{lmodern}
\usepackage{textcomp}
\usepackage{microtype}
\DisableLigatures[-]{encoding = T1, family = tt*}
\usepackage[a4paper,margin=25mm]{geometry}
\usepackage{parskip}
\usepackage{xcolor}
\usepackage{array}
\usepackage{booktabs}
\usepackage{tabularx}
\usepackage{longtable}
\usepackage{enumitem}
\usepackage{caption}
\usepackage{fancyhdr}
\usepackage{upquote}
\usepackage{listings}
\usepackage[hidelinks]{hyperref}

\newcommand{\hdlvetversion}{0.1.0}


\setlist{itemsep=2pt,topsep=4pt}
\setlist[description]{font=\bfseries,leftmargin=1.5em,labelindent=0pt}
\emergencystretch=3em
\raggedright
\setcounter{tocdepth}{1}
\setcounter{secnumdepth}{2}
\raggedbottom

\hypersetup{pdftitle={hdl-vet user manual},pdfauthor={hdl-vet}}

\pagestyle{fancy}
\fancyhf{}
\fancyhead[L]{\small hdl-vet \hdlvetversion}
\fancyhead[R]{\small\nouppercase{\leftmark}}
\fancyfoot[C]{\small\thepage}
\renewcommand{\headrulewidth}{0.3pt}
\fancypagestyle{plain}{\fancyhf{}\fancyfoot[C]{\small\thepage}\renewcommand{\headrulewidth}{0pt}}

% Listings: every example of the manual is a file under examples/, shown with \lstinputlisting.
\definecolor{ruleblue}{RGB}{70,90,140}
\definecolor{notegray}{RGB}{110,110,110}
\lstdefinelanguage{transcript}{morecomment=[l]{\$\ }}
\lstdefinestyle{base}{%
  basicstyle=\ttfamily\footnotesize,
  breaklines=true,
  breakatwhitespace=true,
  postbreak=\mbox{\textcolor{notegray}{\(\hookrightarrow\)}\space},
  keepspaces=true,
  columns=fullflexible,
  upquote=true,
  showstringspaces=false,
  frame=leftline,
  framerule=1.2pt,
  rulecolor=\color{ruleblue},
  xleftmargin=1.2em,
  xrightmargin=0pt,
  aboveskip=0.8em,
  belowskip=0.8em,
  extendedchars=false}
\lstdefinestyle{script}{style=base,language=Python,keywordstyle=\bfseries,commentstyle=\color{notegray}\itshape}
\lstdefinestyle{console}{style=base,language=transcript,commentstyle=\bfseries}
\lstdefinestyle{source}{style=base,commentstyle=\color{notegray}\itshape,morecomment=[l]{//}}
\lstset{style=base}
% The file name stands above each listing and stays with it.
\let\lstinputlistingorig\lstinputlisting
\renewcommand{\lstinputlisting}[2][]{%
  \par
  \ifdim\pagegoal=\maxdimen\else\ifdim\dimexpr\pagegoal-\pagetotal\relax<7\baselineskip\newpage\fi\fi
  \addvspace{0.9em}\noindent
  {\footnotesize\bfseries\ttfamily\detokenize{#2}}\par\nobreak\vskip0.2em
  \lstinputlistingorig[#1,aboveskip=0pt]{#2}}
\captionsetup{font=small,labelfont=bf}

% Names of commands, options and code in running text.
\newcommand{\cmd}[1]{\texttt{\detokenize{#1}}}
\newcommand{\code}[1]{\texttt{\detokenize{#1}}}
\newcommand{\opt}[1]{\texttt{\detokenize{#1}}}
\newcommand{\file}[1]{\texttt{\detokenize{#1}}}
\newcommand{\term}[1]{\emph{#1}}
\newcommand{\apiname}[1]{\texttt{\detokenize{#1}}}

% The scripting API: api.tex has one \apisig{Class.name}{signature} for every public name.
% \sig{Class.name} prints the name and its signature; a property prints without a parameter list.
\newcommand{\apisig}[2]{\expandafter\def\csname apisig@\detokenize{#1}\endcsname{\detokenize{#2}}}
\newcommand{\sig}[1]{%
  \ifcsname apisig@\detokenize{#1}\endcsname
    \begingroup
    \edef\sigtext{\csname apisig@\detokenize{#1}\endcsname}%
    \edef\sigprop{\detokenize{property}}%
    \texttt{\detokenize{#1}}%
    \ifx\sigtext\sigprop\else\texttt{\sigtext}\fi
    \endgroup
  \else
    \PackageError{hdl-vet}{Unknown API name #1}{Run the manual test with HDL_VET_MANUAL_UPDATE=1.}%
  \fi}
% Generated. Checked by tests/test_manual.py against hdl_vet.ops; rewrite it with HDL_VET_MANUAL_UPDATE=1.
\apisig{Instance.allow_changes_below}{()}
\apisig{Instance.connect}{(port, expr, was=None)}
\apisig{Instance.disconnect}{(port, was=None)}
\apisig{Instance.info}{property}
\apisig{Instance.name}{property}
\apisig{Instance.paths}{property}
\apisig{Instance.replace_module}{(module, *, ports=None, params=None, connect=None, set_params=None, drop_ports=(), drop_params=(), unconnected=(), was=None)}
\apisig{Instance.set_param}{(name, value, was=None)}
\apisig{Instance.target}{property}
\apisig{Module.add_assign}{(target, expr)}
\apisig{Module.add_param}{(name, default, type="int")}
\apisig{Module.add_port}{(name, direction, type="logic", width=None)}
\apisig{Module.add_wire}{(name, width=None, type="logic")}
\apisig{Module.info}{property}
\apisig{Module.instance}{(name)}
\apisig{Module.instances}{(module=None, name=None, *, count)}
\apisig{Module.name}{property}
\apisig{Module.paths}{property}
\apisig{Module.remove_assign}{(target, was=None)}
\apisig{Module.remove_instance}{(name, module=None)}
\apisig{Module.remove_param}{(name)}
\apisig{Module.remove_port}{(name)}
\apisig{Module.remove_wire}{(name)}
\apisig{Vet.add_file}{(path)}
\apisig{Vet.check}{(scripts=())}
\apisig{Vet.design}{property}
\apisig{Vet.expect}{(kind, path="*", name="*", *, exact=False)}
\apisig{Vet.finish}{(scripts=())}
\apisig{Vet.group}{()}
\apisig{Vet.module}{(name, instances=None)}
\apisig{Vet.module_paths}{(module)}
\apisig{Vet.open_block}{()}
\apisig{Vet.permit}{(kind, path="*", name="*", *, exact=False)}
\apisig{Vet.propagate_param}{(name, to, *, default=None)}
\apisig{Vet.punch_port}{(name, direction, to, *, signal=None, width=None, type=None, all_instances=False)}
\apisig{Vet.restore}{(path)}
\apisig{Vet.save}{(path, scripts=())}
\apisig{Vet.scope}{(module)}
\apisig{Vet.splice}{(path, start, end, text, old=None)}
\apisig{Vet.text}{(path)}


% One reference entry: the name and signature, then labelled points.
\definecolor{headbg}{RGB}{226,231,244}
\makeatletter
\newenvironment{entry}[1]{%
  \par\addvspace{1.8ex}\noindent
  \colorbox{headbg}{\begin{minipage}{\dimexpr\linewidth-2\fboxsep\relax}\raggedright\bfseries #1\end{minipage}}%
  \par\nobreak\vskip0.4ex
  \@beginparpenalty=10000
  \begin{description}[leftmargin=2.2em,labelindent=0.4em,itemsep=1pt,topsep=2pt]}%
  {\end{description}}
\makeatother

\title{hdl-vet\\[0.4em]\large User manual}
\author{Version \hdlvetversion}
\date{}

\begin{document}

\maketitle
\tableofcontents

\chapter{Introduction}
\label{ch:intro}

\section{What hdl-vet is for}

A generator writes your SystemVerilog. You need a few changes in its output:
a new port on the top module, a tie-off that is no longer a constant, a
technology cell in place of a generic module. If you edit the generated files
by hand, you lose the edits the next time the generator runs.

hdl-vet keeps the edits in \term{scripts}. A script is a short Python file. When the
generator runs again, you run the scripts again, and you get the same edits on
the new files. Each edit is checked. If the result is not what the script
asked for, hdl-vet writes nothing.

hdl-vet uses \texttt{pyslang}, the Python binding of the slang compiler, to
understand the design: its modules, instances, ports, parameters and
connections. It makes each edit as a replacement of a range of bytes in the
source text. After the edits, it loads the edited design again and compares it
with the design you started from. This last step is the \term{sign-off}.

\section{The words used in this manual}

Each word has one meaning in this manual.

\begin{description}
\item[input tree] The folder that holds the design as the generator wrote it. It is
  the deepest folder that contains every source file, every include folder and the
  \emph{base} of the file list (the folder that relative file names in the file list start
  from; chapter~\ref{ch:first}). hdl-vet never changes the input tree. The file list
  itself may lie outside the input tree.
\item[output tree] A new folder that hdl-vet writes. It has the same layout as the
  input tree. It holds the files with your edits, and the files without edits as
  exact copies. A file list that lies outside the input tree is written at the top of
  the output tree.
\item[file and path] A \emph{file} is a file on disk. A \emph{path} is a place in the
  design: a chain of instance names that starts with the top module, for example
  \code{soc_top.u_periph.u_timer0}. This manual uses \emph{file} for the first and
  \emph{path} only for the second.
\item[script] A Python file that defines the function \code{apply(vet)}. It names the
  edits you want.
\item[operation] One edit that hdl-vet knows, for example \apiname{Module.add_port}
  or \apiname{Instance.connect}. A script is a list of calls to operations.
\item[applied] An operation is \term{applied} when it changes the text.
\item[skipped] An operation is \term{skipped} when the design already shows the
  result it asks for. A skipped operation changes nothing. This is why you can
  run a script a second time.
\item[drift] The design is in a state that the operation does not expect: it
  is neither the state before the edit nor the state after it. The operation
  stops with an error. Chapter~\ref{ch:scripts} explains how \code{was=} detects drift.
\item[queued] An operation does not change the text at once. It records the edit in a
  \term{queue}. hdl-vet applies the queue to the text when it next loads the design.
\item[baseline] The design as the compiler first loaded it, before any edit. The
  sign-off compares the final design with the baseline.
\item[elaborated, copy] The compiler builds the design from the top module
  downwards. A module is \term{elaborated} once for each place where it is
  instantiated, and each of these is a \term{copy} of the module. \code{timer} is
  instantiated twice in the example design, so it has two copies. A module that
  nothing instantiates below the top module is not elaborated.
\item[session] The state of a run: the text of every file with the edits so far, the
  queue, the expectations, the checks, the log, the notes, the added files and the
  scripts that ran. You can save a session and continue it later
  (chapter~\ref{ch:session}).
\item[checked operation] An operation that tries its edit on the design before it
  returns, and is refused if the result is wrong (chapter~\ref{ch:scripts}).
\item[expectation] A statement that a structural change must happen
  (\code{expect}) or may happen (\code{permit}) during the run. Operations
  register their expectations themselves. You write one by hand only when you
  edit the text yourself.
\item[sign-off] The final check of a run. hdl-vet loads the edited design and
  compares it with the baseline. If anything is wrong, it writes no output tree.
  Chapter~\ref{ch:signoff} lists everything that it checks.
\end{description}

\section{What you can rely on}

\begin{itemize}
\item The input tree is never changed.
\item The output tree is a copy of the input tree with your edits.
\item Text that no operation edits stays byte-identical. hdl-vet does not
  reformat anything. An operation inserts new text with the line ending of the
  file it edits.
\item Every file that changes gets one extra comment line,
  \code{// Modified by hdl-vet; scripts: ...}, with the names and the hashes
  of the scripts that ran.
\item The sign-off decides whether the output tree is written. If it fails,
  the exit status is 1 and no output tree is written.
\item The same input and the same scripts give the same output, byte for byte.
\end{itemize}

\section{What hdl-vet does not do}

\begin{itemize}
\item It does not check what your logic does. The sign-off compares the
  structure of the design (instances, ports, parameters, connections) and the
  diagnostics of the compiler. It does not simulate. A change of logic that
  you make with \apiname{Vet.splice} passes the sign-off unless you declare it.
  Chapter~\ref{ch:limits} gives an example.
\item It does not edit instances in generate blocks, instance arrays, or
  instances that come from a macro.
\item It sees one configuration of the design: the defines of your file list.
\item It does not write the design from nothing. It edits what a generator
  wrote.
\end{itemize}

Chapter~\ref{ch:limits} collects what is missing and what is known to be wrong.
No other chapter hints at the future.

\section{How to read this manual}

Chapter~\ref{ch:first} takes you through a first project, from start to end.
Read it first. Chapters~\ref{ch:scripts} to~\ref{ch:signoff} explain scripts,
operations and the sign-off. Chapters~\ref{ch:cli} to~\ref{ch:session} describe
the command line, the project file and saved sessions. Chapter~\ref{ch:python}
lists the functions you use when you call hdl-vet from your own Python code.

Every example in this manual is a file, and a test runs it. There are two kinds.
A \term{script} is shown with its file name. A \term{transcript} shows commands
that you type, and what they print. In a transcript, a line that starts with
\texttt{\$ } is a command. The lines below it are its output. A last line
\code{[exit status 2]} shows a non-zero exit status.
A transcript runs in a copy of the folder \file{examples}, which holds the
example design. Where a file name of your own machine would appear, the transcript
shows \code{<examples>}. The program prints its own version as
\code{<version>}.

The example design is small: \code{soc_top} instantiates a reset synchroniser
\code{rst_sync} and a block \code{periph_wrap}, which instantiates the counter
\code{timer} twice. Chapter~\ref{ch:first} shows its files.

\chapter{Installing}
\label{ch:install}

hdl-vet needs Python 3.11 or newer. It needs one other package, \code{pyslang}
version 12.1.0. The installer fetches it for you. The equivalence check needs one more
package. See the section about the extra \code{equiv} below.

You install hdl-vet from PyPI, from a checkout of its repository, or from a
wheel file. Use a virtual environment so that \code{pyslang} does not mix with
other packages.

\section{The systems it runs on}
\label{sec:platforms}

hdl-vet runs on Linux, on macOS and on Windows. CI tests it on all three: on Linux, on
macOS on Apple silicon, and on Windows.

macOS on Intel is not tested. It installs, because the built file that \code{pyslang}
gives for macOS serves both kinds of processor, and nothing in hdl-vet reads the kind.
Nobody runs the tests there, so this manual says what is true and no more.

No compiler is needed on any system. Both packages that hdl-vet needs come as built
files. \code{pyslang} has one for Linux, for macOS and for Windows. The Yosys of the
extra \code{equiv} is one file that holds WebAssembly. It runs anywhere, and it is the
same build on every system, so the equivalence check gives the same answer on all of
them.

On macOS, the Python that comes with the system is older than 3.11. Install a newer one,
for example with Homebrew or with uv, and make the virtual environment from that one.
\cmd{python3 --version} tells you which one you have.

\section{From PyPI}

\begin{enumerate}
\item Create and enter a virtual environment:
  \cmd{python -m venv .venv}, then activate it as your shell requires.
\item Install the package: \cmd{pip install hdl-vet}
\end{enumerate}

The name of the package is \code{hdl-vet}. The name of the command is \cmd{vet}.

\section{From a checkout}

\begin{enumerate}
\item Create and enter a virtual environment:
  \cmd{python -m venv .venv}, then activate it as your shell requires.
\item Go to the folder of the checkout, the one that holds \file{pyproject.toml}.
\item Install it: \cmd{pip install .}
\end{enumerate}

For work on hdl-vet itself, \cmd{pip install -e .} makes an editable install.

\section{From a wheel}

A wheel is one file with the extension \file{.whl}. If you have one,
install it with \cmd{pip install FILE.whl}. To make a wheel from a checkout,
use any build tool that reads \file{pyproject.toml}, for example
\cmd{uv build --wheel}. The wheel appears in the folder \file{dist}.

\section{The extra for the equivalence check}
\label{sec:extra}

\apiname{Module.override_net} changes logic. The sign-off proves the change with Yosys
(section~\ref{sec:equiv}). Yosys comes from the package \code{yowasp-yosys}. hdl-vet pins
one version of it, because a new version of Yosys can change an answer.

The extra \code{equiv} installs this package. Give the name of the extra in square
brackets. Quote the argument, so that your shell does not read the brackets:

\begin{itemize}
\item \cmd{pip install "hdl-vet[equiv]"} installs the release from PyPI.
\item \cmd{pip install ".[equiv]"} installs a checkout.
\end{itemize}

Without the extra, hdl-vet works as before. The one exception is a logic operation. Without
Yosys, \apiname{Module.override_net} refuses the call unless it has a waiver. The message
names \code{hdl-vet[equiv]}.

\section{Check the installation}

Both of these print the program name and the version, for example
\code{hdl-vet 1.2.3}, and nothing else:

\begin{itemize}
\item \cmd{vet --version}
\item \cmd{python -m hdl_vet --version}
\end{itemize}

The commands \cmd{vet} and \cmd{python -m hdl_vet} are the same program. The
transcripts in this manual use \cmd{vet}.

\chapter{A first project}
\label{ch:first}

This chapter takes you from an empty project to a verified output tree. It uses
the example design that comes with the manual.

\section{Before you start}

You can run every command of this chapter yourself. Work in a copy of the example
folder, because one step overwrites a file of the design. The example folder is
\file{docs/manual/examples} in the repository. In a POSIX shell, copy it with
\cmd{cp -r docs/manual/examples /tmp/vet-examples}, then \cmd{cd /tmp/vet-examples}.
In PowerShell, use \texttt{Copy-Item -Recurse docs\textbackslash manual\textbackslash examples
\$env:TEMP\textbackslash vet-examples}, then \texttt{cd \$env:TEMP\textbackslash vet-examples}.
You need \cmd{vet} installed (chapter~\ref{ch:install}).

The transcripts below are one session, shown in parts. A test replays each part in a
fresh copy of the folder, so each part starts with the commands of the parts before it
that it needs (\cmd{vet init}, \cmd{cp}, \cmd{cd}). If you follow along in one copy,
skip the commands that you have already typed. Every part that enters the folder
\file{myproj} leaves it again with \cmd{cd ..}, so you are back in the example folder.

\section{The example design}

The design has four source files and a file list. The file list names the
sources, one per line.

\lstinputlisting[style=source]{examples/design/design.f}

The top module is \code{soc_top}. It has a reset synchroniser and a peripheral
block. The reset synchroniser is a hand-written module.

\lstinputlisting[style=source]{examples/design/rtl/rst_sync.sv}

The peripheral block is generated. It instantiates the counter \code{timer}
twice. \code{timer} is a hand-written module too.

\lstinputlisting[style=source]{examples/design/rtl/periph_wrap.sv}
\lstinputlisting[style=source]{examples/design/rtl/timer.sv}

The top module is generated. It ties the input \code{test_i} of the
synchroniser to the constant \code{1'b0}.

\lstinputlisting[style=source]{examples/design/rtl/soc_top.sv}

Suppose you want a new input port \code{test_en_i} on \code{soc_top}, and you
want it to drive \code{test_i} of the synchroniser instead of the constant. This
is the job of the first script.

\section{The file list and the input tree}

hdl-vet reads the design from a \emph{file list}, a text file of the kind that
simulators read. It understands these elements:

\begin{itemize}
\item a source file, one per line;
\item \code{+incdir+a+b} for include folders, and \code{+define+A=1+B} for defines;
\item \code{-f FILE}, a nested list whose file names start from the base of the top-level
  list, and \code{-F FILE}, a nested list whose file names start from its own folder;
\item comments, from \code{//} to the end of the line, and lines that start with \texttt{\#};
\item environment variables, as \texttt{\$VAR}, \texttt{\$\{VAR\}} or \texttt{\$(VAR)}.
  A variable that is not set is an error.
\end{itemize}

Any other option is an error, so that nothing is dropped without notice.

Relative file names in the file list start from the \term{base}. The base is the folder of
the top-level file list, unless you give \opt{--base}. Use \opt{--base} when your tool
flow writes file lists for the folder where the simulator runs.

The \term{input tree} is the deepest folder that contains every source file, every
include folder and the base. In the example it is the folder \file{design}. The output
tree has the same layout. If the file list lies outside the input tree, hdl-vet writes
it at the top of the output tree, and the tree is not made wider by it. A file list that
sits far from its sources can still make the tree wider than you expect, when the base
is far from the sources too. Then give \opt{--base} with the folder that the sources are
relative to.

\section{Make a project}

A \term{project} is a folder with a project file, \file{hdl-vet.toml}, and your
scripts. The project file names the design, the scripts and the output folder, so
that you can run hdl-vet with no arguments. Chapter~\ref{ch:project} describes it
in full. You make a project with \cmd{vet init}.

\cmd{vet init} needs the file list (\opt{--filelist}, short \opt{-f}) and, if the
design has more than one top-level module, the name of the top module
(\opt{--top}). The example design has one, but we name it anyway. Always name the top
module in your own work. Chapter~\ref{ch:signoff} shows why.

\lstinputlisting[style=console]{examples/console/init.txt}

Read the transcript from the top.

\begin{itemize}
\item \cmd{vet init} loads the design. If the design has errors, \cmd{vet init}
  refuses. Otherwise it writes five files into the new folder \file{myproj}, and
  prints one line for each:
  the project file, a starter script, a \file{README.md}, a \file{.gitignore} and
  a note in the folder \file{cells}, where you keep the Verilog files that your scripts
  add (chapter~\ref{ch:operations}).
\item The project file names the file list relative to itself:
  \code{../design/design.f}.
\item The starter script changes nothing. It names the top module, lists the
  instances of the top module in comments, and shows operations that you can
  copy. It is a valid script: if you run \cmd{vet apply --dry-run} right after
  \cmd{vet init}, the answer is \code{0 applied, 0 skipped}.
\end{itemize}

\section{Write the first script}

A script is a Python file with one function, \code{apply}. hdl-vet calls it with
one argument, the \code{vet} object, which you use to find modules and to call
operations (chapter~\ref{ch:scripts}). Here is the first script.

\lstinputlisting[style=script]{examples/scripts/first.py}

The script does three things.

\begin{enumerate}
\item \code{vet.module("soc_top")} finds the module \code{soc_top} and gives
  you a handle for it.
\item \code{top.add_port("test_en_i", "input")} adds an input port. The type
  defaults to \code{logic}.
\item \code{connect("test_i", "test_en_i", was="1'b0")} connects the port
  \code{test_i} of the instance \code{u_rst_sync} to the new signal. The
  argument \code{was} says what you expect to find there now. If the design shows
  something else, for example after the generator changed, the operation stops
  with an error and changes nothing. This protects you from editing a design
  that is not the one you wrote the script for.
\end{enumerate}

In the next transcripts, the command \cmd{cp} puts this script in place of the starter
script, \file{myproj/scripts/01_example.py}. In your own project you edit that file, or
you add scripts to the list \code{scripts} in the project file.

\section{Look before you write}

The command \cmd{vet apply --dry-run} runs the scripts and the sign-off, and
prints what would change, as a unified diff. It writes no output tree.

\lstinputlisting[style=console]{examples/console/dry_run.txt}

The first line of the output is \code{project: ...}. Whenever hdl-vet uses a project
file, it says so first, so that you always know which project runs. The diff follows.
The last line, \code{2 applied, 0 skipped}, says that two operations changed the text.

Look at the diff. There is one line that you did not ask for. It is the header
line, \code{// Modified by hdl-vet; scripts: ...}. hdl-vet puts it into every
file that it changes, directly after the comment block at the top of the file. It
names each script and the first 12 hexadecimal digits of its hash. Look
also at the new port. It is not aligned with the other ports. hdl-vet never
reformats existing text, and the text it inserts is plain.

\section{Write the output tree}

\cmd{vet apply} without \opt{--dry-run} writes the output tree to the folder that
the project file names (\file{out}).

\lstinputlisting[style=console]{examples/console/run.txt}

The output is the folder, the counts, and one line for each changed file. The folder is
shown with its full name, because it comes from the project file.

If the output folder exists and is not empty, hdl-vet refuses to replace it,
unless you give \opt{--force}. Chapter~\ref{ch:cli} has the details.

\section{Look at the output tree}

The output tree has the same layout as the input tree.

\lstinputlisting[style=console]{examples/console/tree.txt}

\begin{itemize}
\item Files that no operation changed are exact copies.
\item The file list is written again, in \file{design.f}. It has one source per
  line. Nested lists are flattened into it, and it also carries the \code{+incdir+}
  and \code{+define+} lines. Its file names are relative to itself.
\item \file{hdl-vet.json} records what happened: the version, the top module,
  the scripts with their hashes, and for every changed file the hash before and
  after. \code{files} lists every file of the tree except \file{hdl-vet.json}.
  \code{added} lists files that scripts brought in (chapter~\ref{ch:operations}).
  There is no time stamp and no file name of your machine in it.
\end{itemize}

\section{Run it again on a regenerated design}

Next week the generator runs again. It writes a new \file{soc_top.sv}. In the
example, the new file has one more output, \code{busy_o}. You run the same
script on the new files.

\lstinputlisting[style=console]{examples/console/regenerated.txt}

The second command in the transcript, \cmd{cp}, does the work of the generator: it copies
the file \file{regenerated/rtl/soc_top.sv} over the design. This is why you work in a
copy. The same script runs again and gives the same edits, now on the new file. The
output tree has the new port \code{busy_o} and your new port \code{test_en_i}.

Two more things in the transcript.

\begin{itemize}
\item The output tree is not a valid input. Its files carry the header line, and
  hdl-vet refuses to read a file that carries it: \emph{the input must be an
  unmodified tree}. Always start from the generator's files.
\item The last command repeats the third command. The output folder
  \file{out/second} exists, so hdl-vet stops. The message names the option
  \opt{--force}, which replaces the folder.
\end{itemize}

The transcript uses \opt{--no-project}, which tells hdl-vet to ignore a project
file. The example folder has none. It is better to be explicit when you write
commands that other people run.

\chapter{Writing scripts}
\label{ch:scripts}

\section{What a script is}

A script is a Python file. It defines one function, \code{apply}, with one
parameter, the \code{vet} object. hdl-vet calls \code{apply(vet)} once.

\lstinputlisting[style=script]{examples/scripts/first.py}

hdl-vet reads the script as text and runs it. It does not import it, so it
leaves no \file{__pycache__} folder next to it. A script that does not define
\code{apply} is an error. So is a script that you name twice.

Whatever a script prints goes to the standard output of \cmd{vet}. Section~\ref{sec:program}
says what this means for your safety.

A script can fail in three ways. All of them exit with status 2 and write nothing
(chapter~\ref{ch:cli} lists the exit statuses).

\begin{itemize}
\item An operation refuses, or finds drift. hdl-vet prints \code{vet: script NAME:} and the
  reason. The error is a \code{VetError}, or a \code{DriftError}, which is a kind of
  \code{VetError}.
\item The script raises another Python exception. hdl-vet prints
  \code{vet: script NAME raised TYPE: MESSAGE} and then a Python traceback.
\item A raw edit with \apiname{Vet.splice} overlaps an earlier one. The error is a
  \code{SpliceError}. It is not a \code{VetError}, so \code{except VetError} in your script
  does not catch it, and the message is only \code{vet: [2, 7) overlaps [0, 5)}, with the
  two ranges and without the name of the script.
\end{itemize}

\section{Find modules and instances}

You start with \sig{Vet.module}. It takes the name of a module and gives you a
\term{module handle}. You edit a module through its handle. From a module handle
you get \term{instance handles} with \sig{Module.instance} (one instance, by name)
and \sig{Module.instances} (all instances that match a filter).

\lstinputlisting[style=script]{examples/scripts/find.py}

The next two scripts fail on purpose. The first has a spelling mistake in a module
name. The second edits \code{timer} without \code{instances}. In the transcript, look at
the output of \file{find.py}: it prints the paths and the names that the list below
explains. Then look at the two error messages.

\lstinputlisting[style=script]{examples/scripts/typo.py}
\lstinputlisting[style=script]{examples/scripts/needs_count.py}

\lstinputlisting[style=console]{examples/console/find.txt}

\begin{itemize}
\item \apiname{Vet.module} refuses a name that the design does not have, and suggests
  the closest names. It also refuses a module that is not instantiated under the
  top module, and a module that comes from an added file
  (\apiname{Vet.add_file}, chapter~\ref{ch:operations}). hdl-vet edits only the
  modules that it can see in the elaborated design.
\item \apiname{Module.instance} takes the name of one instance in the module. It
  refuses a name that is not there, with suggestions.
\item \apiname{Module.instances} takes an optional \code{module} (the module that is
  instantiated, an exact name), an optional \code{name} (a pattern with
  \code{*} and \code{?}), and \code{count}, which is required: you say how many
  instances you expect. A different number is an error that lists what matched.
  This protects you when a generator changes the number of instances.
\item \sig{Module.paths} is the list of the paths of the module in the
  design, for example \code{soc_top.u_periph.u_timer0}. The first part is the
  name of the top module. \sig{Instance.paths} lists the paths of one instance, one
  for each path of its module. \sig{Vet.module_paths} gives the paths of any
  module by name. These paths describe the design as hdl-vet read it at the start
  of the run. They do not change when your script edits the design.
\item \code{top.path} is the \emph{file} that holds the module. It is not one of the
  paths above. It is an attribute of the model that \code{top.info} returns, and the
  handle forwards it (see the last item). Scripts use it to read and edit the text of the
  module with \apiname{Vet.text} and \apiname{Vet.splice}.
\item \sig{Module.name} and \sig{Instance.name} are the names. \sig{Instance.target}
  is \code{module.instance}, the text that hdl-vet uses in its log and its
  messages.
\item \sig{Module.info} and \sig{Instance.info} give the model of the compiler for
  the module or the instance: ports, instances, connections, and so on. The
  handles forward attribute lookups to it, so \code{top.info.ports} and
  \code{top.ports} are the same list. Read these models; do not change them.
  A handle looks the design up again at each call, so it is never out of date.
\end{itemize}

\subsection*{A module that exists more than once}

\code{timer} is instantiated twice, so \code{periph_wrap} has two instances of
it, and the design has two copies of \code{timer}. When you edit a module, the edit
changes every copy. hdl-vet makes you say that you know.

\begin{itemize}
\item \code{vet.module("timer", instances=2)} confirms that you expect exactly
  two copies. Without \code{instances}, an edit of \code{timer} is refused, with a
  message that lists the paths (the last command of the transcript above). If the
  number is wrong, \apiname{Vet.module} refuses at once.
\item The same holds for instance operations. An instance inside \code{timer}
  would be inside two copies of \code{timer}. You get the handle through
  \code{vet.module("timer", instances=2)}.
\end{itemize}

\section{Applied and skipped}

Every operation looks at the design before it edits.

\begin{enumerate}
\item If the design already shows the result that the operation asks for, the
  operation is \term{skipped}. It queues nothing, and it registers no
  expectations.
\item If the design is in the state that the operation expects before the edit, the
  operation is \term{applied}: it queues the edit and registers what it
  expects of the sign-off.
\item Otherwise there is \term{drift}: the operation stops with an error that
  says what it found.
\end{enumerate}

Because of the first rule you can run a script a second time on a design that
the script has already changed. Nothing happens twice. The summary line,
\code{N applied, M skipped}, counts the operations of the run. A raw text edit
with \apiname{Vet.splice} is not an operation and is not counted.

\lstinputlisting[style=script]{examples/scripts/skipped.py}

In the next transcript, \file{skipped.py} runs after \file{first.py}. It asks
for the port that \file{first.py} has added, so that operation is skipped. The
summary is \code{3 applied, 1 skipped}.

\lstinputlisting[style=console]{examples/console/applied_skipped.txt}

\subsection*{The argument \texttt{was}}

An operation that \emph{replaces} something takes an optional argument \code{was}.
It says what you expect to find. The operations are \apiname{Instance.connect},
\apiname{Instance.set_param}, \apiname{Instance.disconnect}, \apiname{Module.remove_assign}
and \apiname{Instance.replace_module}. \apiname{Module.remove_instance} has the argument
\code{module}, which works in the same way as \code{was}.

\begin{itemize}
\item If the design shows the new value already, the operation is skipped.
  \code{was} is not looked at.
\item If the design shows exactly the value in \code{was}, the operation
  replaces it.
\item If the design shows another value, and \code{was} is given but does not match,
  the operation stops with a drift error. The error quotes the value that it found.
  This holds for all of these operations.
\item If the design shows another value and \code{was} is \emph{missing}, the
  operations differ. \apiname{Instance.connect} and \apiname{Instance.set_param} stop
  with a drift error, because they would overwrite a value that you did not name.
  \apiname{Instance.disconnect}, \apiname{Module.remove_assign},
  \apiname{Instance.replace_module} and \apiname{Module.remove_instance} do not
  check: without \code{was} (or \code{module}) they do their edit on whatever they
  find.
\end{itemize}

The values are compared as text, by the tokens of the language. White space and
comments do not count. You may write a number: \code{was=2} is the text
\code{2}.

The last command of the transcript runs this script. The design does not show
\code{1'b1} at \code{test_i}, so the operation stops with exit status 2, and
nothing is written.

\lstinputlisting[style=script]{examples/scripts/drift.py}

Drift is an error of the script, not a failed sign-off. The exit status is 2.
Give \code{was} always. It turns a silent wrong edit into a message.

\section{Several scripts}

You give scripts to \cmd{vet apply} with \opt{-s}, as many as you like. If the
project file lists scripts, they run first, and the scripts of the command line
run after them. The scripts run in order, in one session. A script sees what the
earlier scripts have done: when a script has finished, hdl-vet loads the edited
text, so the next script finds the new ports, wires and connections. The sign-off
runs once, at the end, over all scripts.

\section{Groups}

Some operations check themselves. They are the \term{checked operations}:
\apiname{Vet.add_file}, \apiname{Instance.replace_module}, \apiname{Module.remove_instance},
\apiname{Module.remove_wire}, \apiname{Module.remove_assign}, \apiname{Module.remove_port} and
\apiname{Module.remove_param}. A checked operation first applies and loads everything
that was queued. Then it makes its edit and loads the design again. It is
\emph{refused}, and everything it did is taken back, when the edited design

\begin{itemize}
\item cannot be loaded, or
\item has a new diagnostic (a new error or a new warning), or
\item leaves a signal or a port \term{dangling}: something reads it, but nothing
  drives it, or
\item leaves a signal partly undriven, for a removal that took away a driver.
\end{itemize}

A refusal is an error of the script (exit status 2). The message starts with the
operation and its target. After a refusal the session is as it was before the
call. If the script catches the error, it can go on.

A removal often needs a second edit to be clean: you remove the driver of a
signal, and you must also remove the signal or its reader. Each step alone is
refused. A \term{group} makes the comparison once, at the end of the block. You open
it with \code{with vet.group():}, which is \sig{Vet.group}. If the comparison fails, or if
an exception leaves the block, everything done in the block is taken back. A group inside
a group is part of the outer group.

\lstinputlisting[style=script]{examples/scripts/no_group.py}
\lstinputlisting[style=script]{examples/scripts/group.py}

The first script removes the reset synchroniser, and does not reconnect
\code{u_periph}. The wire \code{rst_sync_n} is then read but not driven: the
first operation is refused. The second script does all three steps in a group,
and succeeds.

\lstinputlisting[style=console]{examples/console/group.txt}

\subsection*{Notes}

When a checked operation, or a group, leaves something that nobody uses, it is
not an error. hdl-vet adds a line to the list \code{vet.notes}, and the command
prints it as \code{note: ...}. A note starts with the operation and its target,
or with \code{group:} for a group. Notes never make a run fail. The transcript of
\code{retire.py} in chapter~\ref{ch:operations} shows one.

\section{Scopes}

A \term{scope} limits part of a script to one module and what is below it. It
protects you from a mistake that edits another part of the design. You
open a scope with \sig{Vet.scope}, as a \code{with} block.

\lstinputlisting[style=script]{examples/scripts/scope_periph.py}

The scope is the module \code{periph_wrap}. Everything instantiated below it,
here the two timers, is in the scope too. The design is still the whole design,
and the sign-off still judges all of it. A scope adds refusals. It takes no check
away.

\subsection*{What is inside}

\begin{itemize}
\item On entry, hdl-vet applies and loads everything that is queued, as for a
  group. The module must be one that \apiname{Vet.module} accepts, and it must be
  instantiated under the top module. The top module itself is allowed. The scope
  then holds the whole design.
\item An \emph{instance} is inside when it is an instance of the scope module, or
  lies below one, in the design as it is when you enter the block.
\item A \emph{module} is inside when it has at least one instance, and every
  instance of it is inside. A module that is also instantiated elsewhere is not
  inside, even if some of its instances are.
\end{itemize}

In the example, \code{periph_wrap} and \code{timer} are inside the scope
\code{periph_wrap}. \code{soc_top} and \code{rst_sync} are not.

\subsection*{What is refused}

\begin{itemize}
\item An operation of a module handle or an instance handle that would edit the
  text of a module that is not inside is refused when you call it. This holds for
  a handle that you got before the block, too. An instance operation edits the
  module in which the instance is written.
\item \apiname{Module.remove_port} and \apiname{Module.remove_param} also edit every
  module that instantiates the module. They are refused when one of those modules is
  not inside. For the scope module this is the case, unless it is the top module.
\item A refusal names the operation, the module, the scope, and where the module is
  instantiated. The scope is checked before anything else, so an operation that would
  have been skipped is refused too. Nothing is queued. The block goes on, if your script
  catches the error.
\item Reading is free. \apiname{Vet.module}, \apiname{Module.info}, \apiname{Module.paths} and
  \apiname{Module.instance} work for every module.
\item \apiname{Vet.splice} in a scope is refused unless the whole range lies in
  the text of one module that is inside, from the keyword \code{module} to the
  end of \code{endmodule}. It is refused too when the new text holds one of the
  compiler directives \code{`include}, \code{`define}, \code{`undef},
  \code{`undefineall}, \code{`celldefine}, \code{`endcelldefine}, \code{`line},
  \code{`default_nettype}, \code{`unconnected_drive},
  \code{`nounconnected_drive}, \code{`timescale}, \code{`resetall},
  \code{`pragma}, \code{`begin_keywords} or \code{`end_keywords}, outside a
  comment and outside a string. These directives reach text outside the
  module. A use of a macro and the directives \code{`ifdef}, \code{`ifndef},
  \code{`elsif}, \code{`else} and \code{`endif} are allowed.
\item \apiname{Vet.add_file}, \apiname{Vet.expect} and \apiname{Vet.permit} are not
  restricted.
\item \apiname{Vet.punch_port} and \apiname{Vet.propagate_param} start at the scope module.
  The path in \code{to} begins with the name of the scope module, and the new port
  or parameter ends in the list of the scope module itself. Nothing above it is
  edited. In the example, \code{to="periph_wrap.u_timer0"}. A path that starts
  anywhere else is refused, and the scope module alone is refused. Both operations
  are refused when the scope module is instantiated more than once.
\item A scope cannot be opened inside a scope or inside a group. A group may be
  opened inside a scope. \apiname{Vet.check}, \apiname{Vet.save} and \apiname{Vet.finish}
  are refused inside a scope.
\end{itemize}

The scope module's own ports and parameters may change. The parent's
instantiation is outside the scope. So a new input port is not connected in the
parent, and the sign-off fails, until a later operation outside the block
connects it. The example does that after the block: it adds the port to the top
module and connects \code{u_periph}.

\subsection*{What is checked at the end}

When the block ends, hdl-vet loads the design again and compares it with the
design on entry. It raises an error that says \emph{outside the scope}, and takes
back everything done in the block, when

\begin{itemize}
\item any byte outside the text of the inside modules has changed, or
\item the structure has changed at an instance that is not inside, or
\item a module that is inside names, by a hierarchical path, a signal, a task or a
  function of an instance outside the scope, and the design on entry had no such
  reference in that module.
\end{itemize}

The message names one such reference, even if the block added several. The compiler
does not look at a name in code that it does not elaborate, for example an unused
sequence or property, so such a name is not seen. The last point holds however the text
got there: by a splice, or in the
expression of \apiname{Instance.connect} or \apiname{Module.add_assign}. A path that ends
at or below an instance of the scope module is fine. A name that goes through an
interface port of the module is a use of the port, not a reach outside.
If an exception leaves the block, everything done in it is taken back, as for a
group. A script must close every block it opens: if \code{apply} returns with a
scope or a group still open, \cmd{vet} stops with exit status 2 and names the
script.

The rest of the design can be edited again after the block. A scope is not kept in
a saved session (chapter~\ref{ch:session}). Operations in a scope are in
\code{vet.log} as they are outside of one. The next scripts show a second scope,
and two ways to leave one.

\lstinputlisting[style=script]{examples/scripts/scope_reset.py}
\lstinputlisting[style=script]{examples/scripts/scope_leaves.py}
\lstinputlisting[style=script]{examples/scripts/scope_reaches.py}

\lstinputlisting[style=console]{examples/console/scope.txt}

In the transcript, the first command runs two scripts, each with its own scope. It adds a parameter
to \code{rst_sync} inside a scope, and overrides it on the instance outside the
block. The second script writes to \code{soc_top} inside the scope
\code{periph_wrap}, and is refused at the call. The third script names an
instance of \code{rst_sync} from inside \code{periph_wrap}. Without the scope,
the sign-off accepts this reference. In the scope, it reaches outside, so the
block is refused and taken back.

\section{A script is a program}
\label{sec:program}

A script is a Python program. It runs with your rights: it can read and write any file
that you can. hdl-vet does not run it in a sandbox. The sign-off checks the design that
comes out, not what the script did on your machine. \cmd{vet apply} runs the scripts also
with \opt{--dry-run}, and a project file lists scripts that run when you type
\cmd{vet apply}. Run only scripts and projects that you trust.

\chapter{Operations}
\label{ch:operations}

This chapter describes every operation. It groups the operations by what they
change. Each entry has the same points:

\begin{description}
\item[Writes] What the operation puts into the text.
\item[Refused] When it stops with an error and queues nothing. Chapter~\ref{ch:scripts}
  explains drift errors and refusals.
\item[was] What \code{was=} means, for the operations that have it.
\item[Again] What the operation does when you run it a second time on a design that
  already shows its result. It is \emph{skipped}, or it stops with a drift error.
\item[Example] The script and the transcript that show it.
\end{description}

\section{Rules that hold for all operations}

\begin{itemize}
\item An operation on a module handle changes every copy of the module. If the
  design instantiates the module more than once, you must confirm the number with
  \code{vet.module(name, instances=N)}. Without it, hdl-vet refuses the operation.
  The same holds for an operation on an instance handle. In this case, the module that
  holds the instance counts.
\item hdl-vet edits only text that it can edit safely. Chapter~\ref{ch:limits}
  collects these limits.
  \begin{itemize}
  \item hdl-vet refuses an operation on an instance when the instance comes from a
    macro, is written in a generate block, or is an instance array.
  \item It refuses an operation for connections when the instance uses positional or
    \code{.*} connections. It refuses an operation for overrides when the instance uses
    positional parameter overrides.
  \item It refuses an operation when the list that it edits holds a conditional
    compilation directive, for example \code{`ifdef}.
  \item It refuses an operation on a port list or a parameter list when the module
    header has a macro. Operations on ports need an ANSI port list.
  \end{itemize}
\item A file is a file on disk. A path is a place in the hierarchy (chapter~\ref{ch:intro}).
\item A name that does not exist raises an error. The error has up to three
  suggestions, in the form ``did you mean \ldots?''
\item An operation inserts text with the line ending of the file that the text goes
  into. hdl-vet never reformats text that it does not edit.
\item Each applied operation registers \emph{expectations} for the sign-off. These are
  the structural changes that the operation must cause. The entries name them as
  ``Sign-off requires'' and ``Sign-off permits''. Chapter~\ref{ch:signoff} explains
  how the sign-off uses them.
\end{itemize}

\section{All operations at a glance}

For each operation, the table shows three things:

\begin{itemize}
\item which text the operation edits
\item what \code{was} means for the operation, where it has one
\item what happens when you run the operation a second time on a result that already
  exists
\end{itemize}

The later sections give the details.

{\small
\begin{tabular}{@{}>{\raggedright\arraybackslash}p{0.305\textwidth}>{\raggedright\arraybackslash}p{0.2\textwidth}>{\raggedright\arraybackslash}p{0.28\textwidth}>{\raggedright\arraybackslash}p{0.12\textwidth}@{}}
\toprule
Operation & Edits the text of & \code{was} & Second run \\
\midrule
\apiname{Module.add_port} & the module & none & skipped \\
\apiname{Module.add_param} & the module & none & skipped \\
\apiname{Module.add_wire} & the module & none & skipped \\
\apiname{Module.add_assign} & the module & none & skipped \\
\apiname{Module.override_net} & the module & none & skipped, or drift \\
\apiname{Instance.connect} & the module that holds the instance & the expression now there. Without it, a different expression is drift & skipped \\
\apiname{Instance.disconnect} & the module that holds the instance & the expression now there, optional check & skipped \\
\apiname{Instance.set_param} & the module that holds the instance & the override now there. Without it, a different override is drift & skipped \\
\apiname{Vet.punch_port} & every module on the path & none & skipped \\
\apiname{Vet.propagate_param} & every module on the path above the destination & none & skipped \\
\apiname{Vet.add_file} & no source. The output tree gets a copy & none & skipped \\
\bottomrule
\end{tabular}

\medskip
\begin{tabular}{@{}>{\raggedright\arraybackslash}p{0.305\textwidth}>{\raggedright\arraybackslash}p{0.2\textwidth}>{\raggedright\arraybackslash}p{0.28\textwidth}>{\raggedright\arraybackslash}p{0.12\textwidth}@{}}
\toprule
Operation & Edits the text of & \code{was} & Second run \\
\midrule
\apiname{Instance.replace_module} & the module that holds the instance & the module now used, optional check & skipped, or drift \\
\apiname{Module.remove_instance} & the module that holds the instance & \code{module=}, the module now used, optional check & skipped \\
\apiname{Module.remove_wire} & the module & none & skipped \\
\apiname{Module.remove_assign} & the module & the right-hand side now there, optional check & skipped \\
\apiname{Module.remove_port} & the module, and every module that instantiates it & none & skipped \\
\apiname{Module.remove_param} & the module, and every module that instantiates it & none & skipped \\
\apiname{Vet.splice} & the file that you name & none. \code{old=} guards the bytes & none \\
\bottomrule
\end{tabular}
}

\section{Ports and connections}

\begin{entry}{\sig{Module.add_port}}
\item[Writes] The operation writes a port declaration at the end of the ANSI port list.
  It builds the declaration as direction, type, range and name, for example
  \code{input  logic [3:0] name}. The range comes from \code{width}:
  \begin{itemize}
  \item An integer greater than 1 gives \code{[N-1:0]}.
  \item A string \code{W} gives \code{[W-1:0]}.
  \item A string that starts with \code{[} gives that string as the range, as written.
  \item \code{None} or 1 gives no range.
  \end{itemize}
  If the closing parenthesis is alone at the start of its line, the new port
  goes on a new line, indented like the last port. A comma follows the last
  port. Otherwise the text \code{, declaration} follows the last port. An empty
  list gets the declaration between its parentheses.
\item[Refused] hdl-vet refuses the operation in these cases:
  \begin{itemize}
  \item a direction other than \code{input}, \code{output} or \code{inout}
  \item a module with a non-ANSI port list, or with none
  \item a header with a macro
  \item a conditional directive in the list
  \item a name that something else already uses, or a port of that name with a different
    declaration (a drift error)
  \item a name that is not a plain identifier, such as \code{"a b"}, or is a keyword,
    such as \code{"wire"}
  \item a name that begins with a backslash, or is not a \code{str}
  \end{itemize}
  The operation checks the name when you call it. It queues nothing for a name that it
  refuses.
\item[Again] hdl-vet skips the operation when the module has a port of that name with the
  same declaration.
\item[Example] \file{first.py} (chapter~\ref{ch:first}) and \file{ports.py} below.
  The operation does not connect a new input port in the instances of the module.
  Connect it with \apiname{Instance.connect}. If you do not, the sign-off reports a new
  diagnostic.
\end{entry}

\begin{entry}{\sig{Instance.connect}}
\item[Writes] The operation writes a connection \code{.port(expr)} after the last
  connection. It aligns the name field like the others when they all have the same width.
  An instance with no connections gets the connection between its parentheses. If the
  port is connected to nothing (\code{.port()}), the operation fills in the expression.
  If the port is connected to another expression, the operation replaces that
  expression, but only when \code{was} matches. Only the text between the parentheses
  changes.
\item[Refused] hdl-vet refuses the operation in these cases:
  \begin{itemize}
  \item a port that the module does not have
  \item an empty expression (to leave an output open, use \apiname{Instance.disconnect})
  \item the refusals for instances in the rules above
  \end{itemize}
\item[was] The expression that you expect to find now. When the port is connected to
  another expression, and you give no \code{was} or a different one, the operation stops
  with a drift error. The error quotes what it found. The operation fills an open port
  (\code{.port()}) without \code{was}.
\item[Again] hdl-vet skips the operation when the port already has the same expression.
\item[Example] \file{first.py}, and \file{drift.py} in chapter~\ref{ch:scripts}.
\end{entry}

\begin{entry}{\sig{Instance.disconnect}}
\item[Writes] The operation writes an empty connection \code{.port()}. The connection
  keeps the name field.
\item[Refused] The operation refuses an input port, because an open input floats. The
  message names the instance and the port. It gives two other choices:
  \begin{itemize}
  \item tie the input to a value with \apiname{Instance.connect}
  \item remove the port from the module
  \end{itemize}
  The operation refuses an interface port too.
\item[was] The expression that you expect to find now. A different expression is a drift
  error. Without \code{was}, the operation does not check. It opens the port whatever is
  connected.
\item[Again] hdl-vet skips the operation when the port is already open or has no connection.
\item[Example] \file{ports.py}. After the change, \code{count1_o} of the top
  module has no driver. This is allowed, because a top output that nobody reads can float.
  If something reads the signal, the sign-off says \emph{left dangling}
  (chapter~\ref{ch:signoff}).
\end{entry}

\lstinputlisting[style=script]{examples/scripts/ports.py}

\section{Parameters}

\begin{entry}{\sig{Module.add_param}}
\item[Writes] The operation writes \code{parameter type name = default} after the last
  parameter in the parameter port list. If the module has no such list, the operation
  writes \texttt{ \#(parameter type name = default)} directly after the name of the
  module, as in the transcript below. \code{default} can be any value. The operation
  writes it as text.
\item[Refused] hdl-vet refuses the operation in these cases:
  \begin{itemize}
  \item a header with a macro
  \item a conditional directive in the list
  \item a name that another parameter in the body, or something else, already uses (a
    drift error)
  \item a name that is not a plain identifier, such as \code{"a b"}, or is a keyword,
    such as \code{"module"}
  \item a name that begins with a backslash, or is not a \code{str}
  \end{itemize}
  The operation checks the name when you call it. It queues nothing for a name that it
  refuses.
\item[Again] hdl-vet skips the operation when the module has a parameter of that name with
  the same type and default.
\item[Example] \file{params.py}.
\end{entry}

\begin{entry}{\sig{Instance.set_param}}
\item[Writes] The operation writes an override \code{.name(value)} after the last
  override. When the instance has no overrides, it writes \texttt{ \#(.name(value))}
  directly after the module name. It replaces an existing override when \code{was}
  matches.
\item[Refused] hdl-vet refuses the operation in these cases:
  \begin{itemize}
  \item a parameter that the instantiated module does not have, a \code{localparam}, or a
    type parameter
  \item positional overrides
  \item an instance that shares its parameter list with other instances in one statement
  \item the refusals for instances in the rules above
  \end{itemize}
\item[was] The override text that you expect to find now. A different text is a
  drift error.
\item[Again] hdl-vet skips the operation when the override already has the same text.
\item[Example] \file{params.py}. The sign-off compares the text of the override,
  not the resolved value. A change of the parameter can change port widths below
  the instance. Declare these changes with \apiname{Vet.expect} (see
  section~\ref{sec:byhand}).
\end{entry}

\lstinputlisting[style=script]{examples/scripts/params.py}

\section{Wires and assignments}

\begin{entry}{\sig{Module.add_wire}}
\item[Writes] The operation writes a declaration such as \code{logic [3:0] name;} on a
  new line directly below the line that ends the module header. It indents the
  declaration like the first line of the body. If the body is empty, it indents with two
  blanks.
\item[Refused] hdl-vet refuses the operation in these cases:
  \begin{itemize}
  \item a name that is not a plain identifier, or is a keyword, or begins with a
    backslash, or is not a \code{str}
  \item a module header whose last line holds the beginning of the body, or ends inside a
    block comment
  \item a name that something else uses, or a net of that name with another type (a drift
    error)
  \end{itemize}
\item[Again] hdl-vet skips the operation when the module has a net or variable of that name
  with the same type and range.
\item[Example] \file{skipped.py} (chapter~\ref{ch:scripts}).
\end{entry}

\begin{entry}{\sig{Module.add_assign}}
\item[Writes] The operation writes \code{assign target = expr;} on a new line directly
  above the line of \code{endmodule}.
\item[Refused] hdl-vet refuses the operation in these cases:
  \begin{itemize}
  \item a target that is not a net, variable or port of the module (with suggestions)
  \item an input port
  \item a target that already has a driver
  \item a module whose \code{endmodule} does not start its line
  \end{itemize}
  The target must be a plain identifier.
\item[Again] hdl-vet skips the operation when the module has exactly that assignment.
\item[Example] \file{ports.py}.
\end{entry}

\section{Logic}
\label{sec:logic}

The operation in this section changes what the design computes. The sign-off proves that
the change has no effect when you switch it off (chapter~\ref{ch:signoff}).

\begin{entry}{\sig{Module.override_net}}
\item[Writes] The operation puts a multiplexer on a net. The net keeps its name and its
  readers. A new signal, \code{keep}, holds the value that the driver gives. The default
  name of \code{keep} is the name of the net with \code{_pre} at the end. The operation
  makes three changes in the text of the module:
  \begin{itemize}
  \item It declares \code{keep} with the type and the width of the net. The declaration is
    on a new line below the declaration of the net. For a port, it is below the line
    that ends the module header.
  \item It makes the driver of the net drive \code{keep}. The driver is an \code{assign}
    statement, or the connection of an output pin of an instance.
  \item It writes \code{assign net = enable ? value : keep;} on a new line below the
    driver.
  \end{itemize}
  \code{value} is the text of an expression over signals of the module. \code{enable} is
  the name of an input port of the module. The port is one bit wide. Add it first, with
  \apiname{Module.add_port} or \apiname{Vet.punch_port}. When \code{enable} is 0, the net
  has its old value. When \code{enable} is 1, the net has \code{value}.

  The input design must not have the port of \code{enable}. The proof ties the enable to 0
  in the edited module only.

  \code{value} is one expression, and it has the width of the net. A fill literal
  (\code{'0}, \code{'1}, \code{'x} or \code{'z}) takes the width of the net. hdl-vet finds
  the width of the value in the compiler. It does not guess from the text.
\item[Refused] hdl-vet refuses the operation in these cases:
  \begin{itemize}
  \item a net that the module does not have, or an input port
  \item a net with no driver, with more than one driver, or with a driver for a part of
    the net only
  \item a net that an \code{always} block drives
  \item a net that an inout pin of an instance connects
  \item a net that is an unpacked array, or has an initial value
  \item a net that has an override already. This is a drift error. A net gets one override.
  \item a driver in a generate block or in an instance array
  \item a net that has an unpacked dimension or an initial value
  \item an \code{enable} that is not an input port of the module, or is wider than one bit
  \item an \code{enable} that the input design has as a port already
  \item a \code{keep} that is not a plain identifier, or that the module uses already
  \item an empty \code{value}
  \item a \code{value} that is not one expression. Examples are a comma at the top level, an
    unbalanced bracket, a second statement and a comment.
  \item a \code{value} with a line break. A value is one line.
  \item a \code{value} that reads the net itself, or a name the net drives inside the
    module. The message says \emph{loop}. hdl-vet reads the names from the syntax of the
    value, not from its text. It judges the value by the name each path starts with. A
    route through a child instance, a hierarchical name or a function call is not seen
    (chapter~\ref{ch:limits}).
  \item a \code{value} with a width other than the width of the net. The message says
    \emph{width}.
  \item a module that the design instantiates more than once, when you did not give the
    count (see the rules above)
  \item a \code{waiver} that is not a text, or has a line break
  \item a call without a waiver when Yosys is not installed. The message names
    \code{hdl-vet[equiv]}.
  \end{itemize}
  The operation checks all this when you call it. It queues nothing for a call that it
  refuses. Like the removals, it loads the edited design. If the design has a new
  diagnostic, the operation takes its edits back and stops.
\item[Sign-off requires] When an instance pin drives the net, the sign-off requires a
  \code{connection-changed} of that pin at every path of the module. An \code{assign}
  driver needs no change of the structure.
\item[Waiver] \code{waiver} is a text that says why the proof is not needed, for example
  the name of a test that covers the change. hdl-vet prints it in the report and writes it
  to \file{hdl-vet.json}. With Yosys, the check still runs and the report shows what it said. Only the waiver
  decides whether the sign-off passes. Without Yosys, the operation needs a waiver.
  Chapter~\ref{ch:signoff} gives the rules.

  The waiver is one line. The report shows it in quotes, as JSON does. The waiver of the
  latest call counts. A call that is skipped sets the waiver of that
  override to what the call gives. A call without a waiver takes the waiver away.
\item[The way of the enable] The proof ties the enable to 0. The sign-off also follows the
  enable up to the top module (section~\ref{sec:enable}). It looks at the pin of the enable
  on each instance of the module in every elaboration of the edited design. A pin is
  \emph{off} when it connects a literal with the value 0. A pin is \emph{carried} when the
  compiler resolves its connection to a new input port of the parent. Any other connection is \emph{not off}, and the sign-off fails unless every
  operation that uses the enable has a waiver.
\item[was] The operation has no \code{was}.
\item[Again] hdl-vet skips the operation when the net is driven by the line
  \code{assign net = enable ? value : keep;} that the operation writes. If that line has
  another enable, another value or another kept signal, the operation stops with a drift
  error.
\item[Example] \file{override.py} bypasses the reset of the example design. The script
  \file{override_again.py} runs after it and is skipped. The script
  \file{override_waiver.py} gives a waiver. The script \file{override_refused.py} names an
  enable with two bits. The script \file{override_wrong.py} swaps the two branches of the
  multiplexer by hand. The sign-off then says that the module is not equivalent.

  The script \file{override_levels.py} puts one enable on an override in \code{rst_sync} and
  on an override in \code{soc_top}. The script \file{override_not_off.py} connects the pin of
  the enable of \code{rst_sync} to 1. The sign-off then says that the enable is not off.
\end{entry}

\lstinputlisting[style=script]{examples/scripts/override.py}
\lstinputlisting[style=script]{examples/scripts/override_again.py}
\lstinputlisting[style=script]{examples/scripts/override_waiver.py}
\lstinputlisting[style=script]{examples/scripts/override_refused.py}
\lstinputlisting[style=script]{examples/scripts/override_wrong.py}
\lstinputlisting[style=script]{examples/scripts/override_levels.py}
\lstinputlisting[style=script]{examples/scripts/override_not_off.py}

\lstinputlisting[style=console]{examples/console/override.txt}

The first command shows the diff and the line of the check. The second command shows the
second run. The third command shows a waiver. The fourth command shows a refusal. The fifth
command shows a failed check. The sixth command shows two levels, and the seventh command
shows an enable that is not off.

\section{Through the hierarchy}

A generator makes the hierarchy. A change deep inside the hierarchy means a port or a
parameter in every module on the way. The next two operations do all of this.

The two operations take a hierarchical path \code{to}. The path starts with the name of
the top module, for example \code{soc_top.u_periph.u_timer0}. The modules on the path are
$M_0$ (the top module) to $M_n$ (the destination). hdl-vet checks the whole path
before it changes anything.

\begin{entry}{\sig{Vet.punch_port}}
\item[Writes] For an \emph{input}, a signal enters at the top module and arrives at
  $M_n$.

  $M_0$ gets the input port \code{name}, unless you give \code{signal}. \code{signal} is
  the name of an existing port or net of $M_0$ that is the source. In this case $M_0$
  gets no new port. Each of $M_1\ldots M_n$ gets the input port \code{name}. The
  operation connects the instance of each of these modules in the module above to
  \code{name}. If you gave \code{signal}, it connects the instance in $M_0$ to
  \code{signal}.

  For an \emph{output}, \code{signal} (default \code{name}) is a port or net of $M_n$.
  $M_n$ gets the output port \code{name} and \code{assign name = signal;}. Each of
  $M_{n-1}\ldots M_0$ gets the output port. The operation connects the instance below to
  it.

  If you give \code{type} or \code{width}, the operation declares the port as
  \apiname{Module.add_port} does. If you give neither and there is a signal, the port
  takes the type of that signal as written (an output always has one). Without a signal
  the input is a 1-bit \code{logic}.

  The type must mean the same in every module on the path. For this, you must write
  every identifier in the type as \code{pkg::NAME}. The operation refuses a type that
  uses a parameter or typedef of one module. In this case, give \code{width=}.
\item[Refused] hdl-vet refuses the operation in these cases:
  \begin{itemize}
  \item a path that is the top alone, does not start at the top module, or names an
    instance that does not exist
  \item a segment with an index, instances in generate blocks and instance arrays
  \item an interface port
  \item a signal with unpacked dimensions
  \item a module on the path that the design instantiates more than once
  \end{itemize}
  For such a module, the operation refuses an input unless you give
  \code{all_instances=True}. It always refuses an output, because several copies then
  drive one net. With \code{all_instances=True}, the operation connects every instance of
  each module on the path. Each of these instances must be in a module of the path.
\item[Again] A second identical call skips every piece. If a path is partly
  punched, the operation completes it. If the output signal has the same name as an
  existing output port of $M_n$, the operation does not edit $M_n$.
\item[Example] \file{punch.py} and \file{punch_refused.py}.
\end{entry}

\lstinputlisting[style=script]{examples/scripts/punch.py}

\lstinputlisting[style=script]{examples/scripts/punch_refused.py}

\lstinputlisting[style=console]{examples/console/hierarchy.txt}

In the transcript, \file{punch_refused.py} catches the two refusals and prints
them. The design instantiates the module \code{timer} twice. The input needs
\code{all_instances=True}. The output cannot pass through two copies at all.

\file{punch.py} then does what is possible. It brings the signal
\code{rst_no} out of \code{u_rst_sync} to a new output of the top module. It
adds an input to every copy of \code{timer} down from the top. In the listing of
\file{soc_top.sv}, look at the new lines:

\begin{itemize}
\item the output port \code{sync_rst_o} and the input port \code{tick_i} at the end of
  the port list
\item the connection \code{.sync_rst_o(sync_rst_o)} on \code{u_rst_sync}
\item the connection \code{.tick_i(tick_i)} on \code{u_periph}
\end{itemize}

\begin{entry}{\sig{Vet.propagate_param}}
\item[Writes] The operation makes a parameter of $M_n$ settable from the top, without a
  change of its value. Each of $M_0\ldots M_{n-1}$ gets
  \code{parameter type name = default}. The operation overrides the instance below with
  \code{.name(name)}. The default is the first of these that exists:
  \begin{enumerate}
  \item \code{default}, if you give it
  \item the value that overrides the destination instance now
  \item the default that $M_n$ declares
  \end{enumerate}
  The type is the type of $M_n$. An implicit type stays implicit.
\item[Refused] hdl-vet refuses the operation in these cases:
  \begin{itemize}
  \item a parameter that $M_n$ does not have, a \code{localparam}, a type
    parameter, or a parameter declared in the body
  \item a type or a value that is not the same in every module on the path (the rules of
    \apiname{Vet.punch_port}). In this case, give \code{default=}.
  \end{itemize}
  The design must instantiate each of $M_0\ldots M_{n-1}$ once. The design can
  instantiate $M_n$ itself more than once. The operation does not edit it.
\item[Again] hdl-vet skips the operation when the path already passes the parameter down.
\item[Example] \file{propagate.py}. The value in use must not change. If it changes,
  the sign-off says so. The sign-off does not permit a parameter change at the
  destination, and a check compares the resolved value before and after.
\end{entry}

\lstinputlisting[style=script]{examples/scripts/propagate.py}

\lstinputlisting[style=console]{examples/console/ports.txt}

In the transcript, look at three diffs:

\begin{itemize}
\item \file{ports.py} adds the port \code{status_o} with its assignment, and opens
  \code{count1_o} on \code{u_periph}.
\item \file{params.py} turns the module header into
  \code{module soc_top }\texttt{\#}\code{(parameter int DEPTH = 2) (} and uses the
  parameter in the override.
\item \file{propagate.py} adds a parameter \code{WIDTH} to \code{periph_wrap} and to
  \code{soc_top}, and passes it down. In \code{soc_top} the instance becomes
  \code{periph_wrap }\texttt{\#}\code{(.WIDTH(WIDTH)) u_periph}. In \code{periph_wrap}
  the first timer gets \code{.WIDTH(WIDTH)}. The value stays 8 everywhere.
\end{itemize}

\section{Replacing a module}

You have a module that is not in your design, a technology cell. You want an
instance to use it. The cell is in a file that is not in the file list of the
generator. Two operations do the job. \apiname{Vet.add_file} adds the file to the
design. \apiname{Instance.replace_module} changes the module of the instance.

\lstinputlisting[style=source]{examples/cells/tech_rst_sync.sv}

\begin{entry}{\sig{Vet.add_file}}
\item[Writes] The operation writes nothing in the sources. The file joins the design as
  its last source. hdl-vet never edits it. These things happen in the output tree:
  \begin{itemize}
  \item The output tree has a copy of the file at \file{hdl_vet_added/NAME}.
  \item The file list gets one more line for the file, after the files that belong to the
    design.
  \item \file{hdl-vet.json} lists the file under \code{added}.
  \end{itemize}
  On the console, hdl-vet prints \code{added hdl_vet_added/NAME}. A relative file name
  in the argument \code{path} starts from the folder of the script. It starts from the
  current folder when no file of that name is next to the script. An absolute file name
  stays as it is. Outside a script, a relative file name starts from the current folder.
\item[Refused] hdl-vet refuses the operation in these cases:
  \begin{itemize}
  \item a relative file name that is next to the script and also in the current folder, but
    not one file: give an absolute file name
  \item a file that cannot be read
  \item a second added file with the same name
  \item an input tree that has the folder \file{hdl_vet_added} at its root
  \item a file that includes a header that is not a file of the design
  \item a file whose addition gives an error: a module that exists already, or a syntax
    error
  \end{itemize}
  hdl-vet does not count a warning as new if the file had it when the script added the
  file. hdl-vet counts a warning that appears in the file later, for example after an
  edit elsewhere changes how the file is read.
\item[Again] hdl-vet skips the operation when the file is already part of the design.
\item[Sign-off] The sign-off fails with \code{failed check: add_file NAME: none of
  its modules is instantiated} when no instance, other than a top module, uses a module
  from the file. The last transcript below shows this failure.
\item[Example] \file{swap_cell.py}, \file{add_cell.py}.
\end{entry}

\begin{entry}{\sig{Instance.replace_module}}
\item[Writes] The operation writes the name of the new module in place of the old name.
  It also writes the changes to the connections and overrides that you name. The steps run
  in this order:
  \begin{enumerate}
  \item The operation replaces the module name.
  \item It gives a new name to renamed connections and overrides. The blanks before the
    opening parenthesis stay.
  \item It removes dropped connections and overrides. If an override list becomes empty,
    it removes the whole list.
  \item It adds \code{set_params}, then \code{connect}, then \code{unconnected}. It adds
    the entries of each in the order given.
  \end{enumerate}
  The operation does not reformat existing text.
\item[Rules] Each connection and each override of the instance, and each port
  of the new module, must be covered by one of these rules:
  \begin{itemize}
  \item \code{ports={"old": "new"}} renames a connection. \code{drop_ports=["old"]}
    removes it. The operation keeps a connection that you name in neither, if the new
    module has a port with that name.
  \item \code{params} and \code{drop_params} do the same for parameter overrides.
    This holds for an override of a type parameter too. The operation keeps it when the
    new module has a type parameter of that name. \code{params} renames it onto a type
    parameter. The text of the override does not change.
  \item Every port of the new module needs a connection. The connection is one of these:
    \begin{itemize}
    \item the target of a kept or renamed connection
    \item \code{connect={"new": expr}}
    \item \code{unconnected=["new"]}
    \end{itemize}
    \code{unconnected} is for outputs and inouts. It writes an empty connection.
  \item \code{set_params={"name": value}} adds an override, or replaces a kept one.
    Parameters of the new module that nothing overrides keep their defaults.
  \end{itemize}
\item[Refused] hdl-vet checks the maps first. It reports all problems in one
  message. The problems are:
  \begin{itemize}
  \item something not covered by a rule (the message says \emph{is not accounted for} and
    names each one)
  \item a name that the instance or the new module does not have
  \item an input in \code{unconnected}
  \item two things on one port or parameter
  \item a \code{set_params} name that is not an overridable value parameter
  \item a \code{params} map from a type parameter to a value parameter, or the reverse
  \end{itemize}
  The operation also refuses a module with a non-ANSI port list, an interface port, and
  positional overrides. It tries the edit. It refuses the edit when a kept or renamed port
  has another direction, width or resolved type in the new module. A typedef or a type
  parameter with the same resolved type is no difference. The message names both ports.
  It also refuses the edit when the edit gives another value to a parameter that both modules
  have and nothing overrides.
\item[accept\_types] A list of ports of the new module. A renamed port keeps the name it
  lands on. The list waives a difference of type for these ports only. A difference of
  direction or bit width is never waived. Each waived port adds a note with both types to the
  sign-off report. The operation refuses a name that is not a port of the new module, and a
  waiver that covers nothing.
\item[was] The name of the module that you expect the instance to use now. If the
  instance uses neither the module in \code{was} nor the new module, the operation
  stops with a drift error. The error names the module that the instance uses. Without
  \code{was} the operation does not check. It replaces whatever module the instance uses.
\item[Again] When the instance already uses the new module, hdl-vet compares only
  \code{connect} and \code{set_params}. If every entry is there, hdl-vet skips the operation.
  Otherwise it stops with a drift error that names the first entry that differs or is
  missing.
\item[Sign-off] The sign-off requires \code{module-changed} at every path of the
  instance. It permits every change at these paths and below them. At the end, a check
  confirms that the instance uses the new module.
\item[Example] \file{swap_cell.py}, \file{swap_cell_drop.py} and
  \file{swap_wrong.py}.
\end{entry}

\lstinputlisting[style=script]{examples/scripts/swap_cell.py}

\lstinputlisting[style=script]{examples/scripts/swap_cell_drop.py}

\lstinputlisting[style=script]{examples/scripts/swap_wrong.py}

\lstinputlisting[style=script]{examples/scripts/add_cell.py}

\lstinputlisting[style=console]{examples/console/replace.txt}

In the first script, the ports \code{clk_i}, \code{rst_ni}, \code{test_i} and
\code{rst_no} get the names that the cell uses, and the parameter \code{STAGES}
becomes \code{DEPTH}. The cell has an extra output, \code{SO}. The script
leaves it open with \code{unconnected}.

The second script does the same in another way. It drops \code{test_i}, and it ties
\code{SE} of the cell to a constant with \code{connect}. It also drops the parameter
\code{STAGES} and sets \code{DEPTH}.

In the transcript, you see three things:

\begin{itemize}
\item the diff of the first script
\item the output tree of the second script
\item two failures: \file{swap_wrong.py} does not cover \code{SO} and
  \code{STAGES}. \file{add_cell.py} adds the cell and does not use it.
\end{itemize}

Look at the diff of the first script. The renamed connections, such as \code{.CK  (clk_i)},
keep the blanks that the old names had, so they are not aligned with each other any more.
This is as expected. hdl-vet never reformats bytes that it does not edit, and a rename
changes only the name. The line \code{added hdl_vet_added/tech_rst_sync.sv} after the
diff says that the cell goes into the output tree.

\section{Removing things}

The removal operations are checked operations (chapter~\ref{ch:scripts}). hdl-vet tries
each one on the design. If the operation leaves something dangling, gives a new
diagnostic, or leaves a signal partly undriven, hdl-vet refuses it. It takes back
everything that the operation did. To remove a driver and its reader together, use a group.

The removal takes the lines that it empties. This is the \emph{line rule}. If only
blanks stay on the lines that the removed text occupied, the removal takes these lines
too, with one line break. For an item of a list (a port, parameter, connection or
override), the rule is:

\begin{itemize}
\item If the item is not the last, the removal takes the text from the start of the item
  to the start of the next item.
\item If the item is the last of several, the removal takes the comma after the previous
  item, with the item.
\item If the item is the only item, the removal takes the item. If a list of overrides
  becomes empty, the removal takes the whole list. If a parameter port list becomes empty
  and only blanks stand between the name and \texttt{\#}, the removal takes the whole
  list. This is the text from the end of the module name through the closing parenthesis.
\end{itemize}

\begin{entry}{\sig{Module.remove_instance}}
\item[Writes] The operation removes the instantiation statement (line rule).
\item[Refused] hdl-vet refuses the operation in these cases:
  \begin{itemize}
  \item an instance that comes from a macro, a generate block or an array
  \item a statement that declares several instances
  \end{itemize}
\item[was] \code{module} works as \code{was}. If you give it and the instance uses
  another module, the operation stops with a drift error. Without it, the operation does
  not check which module the instance uses.
\item[Again] hdl-vet skips the operation when the module has no instance of that name.
\item[Sign-off] The sign-off requires \code{instance-removed} at every path. It permits
  any change below these paths.
\item[Example] \file{group.py}.
\end{entry}

\begin{entry}{\sig{Module.remove_wire}}
\item[Writes] The operation removes the declaration of a module-level net or variable
  (line rule).
\item[Refused] hdl-vet refuses the operation in these cases:
  \begin{itemize}
  \item a name that is a port or something other than a net or variable
  \item a declaration that declares several names, comes from a macro or an include, or
    is in a generate block
  \end{itemize}
\item[Again] hdl-vet skips the operation when nothing of that name is declared.
\item[Example] \file{group.py}.
\end{entry}

\begin{entry}{\sig{Module.remove_assign}}
\item[Writes] The operation removes the continuous assignment to \code{target} (line
  rule).
\item[Refused] hdl-vet refuses the operation in these cases:
  \begin{itemize}
  \item several assignments to the target
  \item a statement that holds several assignments, comes from a macro or an include, or
    is in a generate block
  \end{itemize}
\item[was] The right-hand side that you expect. A different one is a drift error.
  Without \code{was}, the operation does not check.
\item[Again] hdl-vet skips the operation when there is no assignment to the target.
\item[Example] \file{retire.py}.
\end{entry}

\begin{entry}{\sig{Module.remove_port}}
\item[Writes] The operation removes the port from the ANSI port list, and the connection
  to it at every instance of the module.
\item[Refused] hdl-vet refuses the operation in these cases:
  \begin{itemize}
  \item a port whose successor takes its direction or type from it (a declaration like
    \code{output logic [3:0] y0, y1})
  \item a non-ANSI port list
  \item a header with a macro
  \item any instance whose connection list cannot be edited
  \end{itemize}
\item[Again] hdl-vet skips the operation when the module has no such port.
\item[Sign-off] The sign-off requires \code{port-removed} at every path when the input
  design had the port. It permits \code{connection-changed} for the port. When
  \apiname{Module.add_port} of this run added the port, the removal forgets that
  expectation instead, so the pair leaves nothing behind.
\item[Example] \file{retire.py}.
\end{entry}

\begin{entry}{\sig{Module.remove_param}}
\item[Writes] The operation removes the parameter from the parameter port list, and its
  override at every instance.
\item[Refused] hdl-vet refuses the operation in these cases:
  \begin{itemize}
  \item a parameter that is declared in the body
  \item a declaration of several names
  \item a header with a macro
  \end{itemize}
\item[Again] hdl-vet skips the operation when the module has no such parameter.
\item[Sign-off] The sign-off requires \code{param-changed} at every path.
\item[Example] \file{retire.py}.
\end{entry}

\lstinputlisting[style=script]{examples/scripts/retire.py}

\lstinputlisting[style=script]{examples/scripts/remove_timer.py}

\lstinputlisting[style=console]{examples/console/remove.txt}

\file{retire.py} removes the top-level output \code{irq_o} from the design.
It removes the assignment and the port in one group. The signal \code{irq} stays.
Nothing reads it any more, so hdl-vet adds a note. The same script removes the parameter
\code{ID}, which nobody uses. The empty parameter list goes too.

hdl-vet refuses \file{remove_timer.py}. If the script removes the second timer, bit 1 of
\code{irq_o} of \code{periph_wrap} has no driver, and something reads it.

\section{Expectations by hand, and raw edits}
\label{sec:byhand}

Sometimes no operation does what you need. Then you edit the text yourself, and you
tell the sign-off what to expect. No operation checks raw edits. The
sign-off checks them, but it checks only what it can see. It can see the structure of
the design and the diagnostics (chapter~\ref{ch:signoff}).

\begin{entry}{\sig{Vet.splice}}
\item[Writes] The operation replaces the bytes from \code{start} up to \code{end} of the
  file \code{path} with \code{text} (a \code{bytes} value). Here \code{path} is a
  \emph{file}. The offsets are offsets into \code{vet.text(path)}. The file is a file of
  the design, for example the attribute \code{path} of a module handle.
\item[Refused] hdl-vet refuses the operation in these cases:
  \begin{itemize}
  \item a file that is not a file of the design, or is not editable (an added file is
    not editable)
  \item a range that overlaps an earlier splice
  \item a range outside the file
  \item an offset that is not an \code{int}, a negative offset, or a \code{start} after
    \code{end}
  \item a \code{text} or an \code{old} that is not \code{bytes}
  \item a range whose bytes differ from \code{old}, if you gave it
  \end{itemize}
  Always give \code{old}. Then the splice cannot go wrong when the text changes. Every
  refusal is a \code{VetError}. Its message starts with \code{splice} and the name of the
  file (chapter~\ref{ch:scripts}). The operation accepts the same range with the same text a
  second time, without a message.
\item[Again] A splice has no skip rule. With \code{old}, a second splice on text that
  is already edited is refused, because the bytes differ.
\item[Example] \file{by_hand.py}. After a splice, the next operation loads the design
  again, because hdl-vet does not know the effect of a splice.
\end{entry}

\begin{entry}{\sig{Vet.text}}
\item[Writes] The operation writes nothing. It returns the text of a file of the design,
  as \code{bytes}, as of the last load. The text does not have edits that are still in the
  queue, such as a raw splice, until the next load. A file that is not a file of the design
  is a \code{VetError}. So is a \code{path} that is neither a \code{str} nor a \code{Path}.
  Its message starts with \code{text} and the name of the file.
\end{entry}

\begin{entry}{\sig{Vet.expect}}
\item[Writes] The operation writes nothing in the text. It declares a change that the
  sign-off must see. \code{kind}, \code{path} and \code{name} are patterns (\code{*},
  \code{?} and \code{[...]}). Capital and small letters differ. \code{kind} is one of
  these:
  \begin{itemize}
  \item \code{instance-added}
  \item \code{instance-removed}
  \item \code{module-changed}
  \item \code{param-changed}
  \item \code{port-added}
  \item \code{port-removed}
  \item \code{port-changed}
  \item \code{port-type-changed}
  \item \code{connection-changed}
  \end{itemize}
  \code{path} is the hierarchical path of the instance. \code{name} is the port or
  parameter. It is empty for the kinds that name none. With \code{exact=True},
  \code{path} is a plain path and not a pattern.
\item[Rules] A change that no expectation or permission covers makes the sign-off fail
  (\emph{unexpected change}). An expectation that matches no change also makes the
  sign-off fail (\emph{expected change not seen}).
\item[Example] \file{by_hand.py}, \file{unexpected.py}, \file{unseen.py}.
\end{entry}

\begin{entry}{\sig{Vet.permit}}
\item[Writes] The operation writes nothing in the text. It works as
  \apiname{Vet.expect}, but the change \emph{can} happen. If no change happens, this is
  not an error. Use it for changes that follow from yours but that you do not have to see.
\item[Example] \file{by_hand.py}.
\end{entry}

\begin{entry}{\sig{Instance.allow_changes_below}}
\item[Writes] The operation writes nothing in the text. It permits every structural
  change strictly below each path of the instance.
\item[Example] \file{below.py}.
\end{entry}

\lstinputlisting[style=script]{examples/scripts/by_hand.py}

\lstinputlisting[style=script]{examples/scripts/below.py}

\lstinputlisting[style=script]{examples/scripts/unexpected.py}

\lstinputlisting[style=script]{examples/scripts/unseen.py}

\lstinputlisting[style=console]{examples/console/by_hand.txt}

\file{by_hand.py} makes two raw edits. It changes the tie-off of \code{test_i}, and
declares this with \code{expect}. It changes the parameter \code{STAGES}, and says
with \code{permit} that this can happen. \file{below.py} edits the enable of the
first timer. The change is below \code{u_periph}, and \code{allow_changes_below}
covers it.

The last two scripts show the two failures of the sign-off:

\begin{itemize}
\item a change that nobody declared
\item a declaration without a change
\end{itemize}

\chapter{Sign-off}
\label{ch:signoff}

The \term{sign-off} is the last step of a run. It decides whether hdl-vet writes the
output tree. You see it when it fails, so this chapter is about what each failure
says and what to do about it.

\section{What the sign-off does}

When all scripts have run, hdl-vet does these steps.

\begin{enumerate}
\item It applies the queued edits to the text.
\item It adds the header line \code{// Modified by hdl-vet; scripts: ...} to every file
  that differs from its original. The header is part of the text that is checked.
  What is checked is byte for byte what would be written.
\item If a file differs, it loads the edited design and compares it with the design
  that you started from, the \term{baseline}.
\item It evaluates the checks that the operations registered.
\end{enumerate}

If every step passes, the output tree is written. If a file did not change, there is
nothing to compare, and the checks run on the baseline. Everything is compared
against the baseline, not against zero: a warning that the generator's design
already has is not your problem, and does not fail the sign-off.

A failed sign-off prints \code{hdl-vet sign-off: FAILED} and one line for each
problem, to the error output. The exit status is 1. No output tree is written. This
holds with \opt{--dry-run} too.

The exit status is 1 when the sign-off fails. Chapter~\ref{ch:cli} lists all exit statuses.

\section{Every line of the report}

Each problem is one line. The lines come in this order. The transcript after the list
shows one failure of each kind.

\begin{description}
\item[the edited design cannot be loaded: REASON] The compiler cannot read the
  edited text, or cannot find the top module. The reason is its message. Look at
  your raw edits, or at an operation that changed a name that the top module needs.
  An operation that fails in this way alone is usually refused before.

\item[new file read by the edited design, not in the baseline: PATH] The edited
  design reads a file that the original design did not read. A raw edit added an
  \code{`include}. Use \apiname{Vet.add_file} to bring in a file on purpose.

\item[new diagnostic: SEVERITY CODE at FILE:LINE: MESSAGE] The compiler reports an
  error or a warning that the baseline did not have. \code{CODE} is the compiler's
  name for it, for example \code{PortWidthExpand}. \code{FILE} and \code{LINE} are
  in the edited text. The compiler reports its errors and its default warnings,
  and hdl-vet switches on five more that an edit of ports can cause:
  \code{ImplicitNet}, \code{PortWidthTruncate}, \code{PortWidthExpand},
  \code{WidthTruncate} and \code{EmptyInputPortConn}. Notes of the compiler are
  not counted. A diagnostic that is in the
  baseline and has only moved because of an edit is not new. Warnings that a file
  already had when a script added it are not new, but a later warning located in that file
  counts. Read the message and the line, and fix the
  script.

\item[left dangling: MESSAGE in module MODULE (FILE:LINE)] Something is read but
  nothing drives it. See the next section.

\item[unexpected change: KIND PATH NAME: BEFORE -> AFTER] The structure of the
  design changed in a way that no operation and no \code{expect} or \code{permit}
  explains. \code{KIND} is the kind of change (see \apiname{Vet.expect}), \code{PATH}
  the instance, \code{NAME} the port or parameter, if any. The values are the text
  before and after. A change that you made by hand needs an \code{expect}. A change
  that you did not expect from an operation shows that the operation touched more than
  you thought.

\item[expected change not seen: KIND path=PATH name=NAME] Something was expected
  (\code{expect}, or an operation that must change the structure) and the design
  does not show it. Another edit may have undone it, or the \code{expect} is wrong.

\item[failed check: MESSAGE] A post-condition failed. See the section after next.
\end{description}

These scripts produce the failures in the transcript. \file{cannot_load.py} renames the
top module with a raw edit. \file{new_file.py} inserts an \code{`include}.
\file{wide_port.py} connects a 1-bit expression to a 2-bit port. \file{dangling.py}
opens an output that something reads. \file{unexpected.py}, \file{unseen.py} and
\file{add_cell.py} are in chapter~\ref{ch:operations}.

\lstinputlisting[style=script]{examples/scripts/cannot_load.py}
\lstinputlisting[style=script]{examples/scripts/new_file.py}
\lstinputlisting[style=script]{examples/scripts/wide_port.py}
\lstinputlisting[style=script]{examples/scripts/dangling.py}

\lstinputlisting[style=console]{examples/console/signoff.txt}

The last-but-one command runs two scripts that fail in different ways. The report has
one line for each problem.

\section{Dangling: an output may float, an input may not}

The sign-off compares the \term{dangling set} of the final design with that of the
baseline. Only new entries fail. The set holds:

\begin{itemize}
\item a net or variable that is read but never driven;
\item an output port that nothing inside its module drives, \emph{but only when
  something reads it}: some instance of the module connects it to a non-empty
  expression, or the module itself reads its signal.
\end{itemize}

An output that nobody reads may float. It is allowed. This is the case when
every instance leaves it open (\code{.p()}), and for an output of the top module.
An input may not float. A new input port that no instance connects is a new
diagnostic. An input connected to nothing with \code{.p()} is a new diagnostic too.
Entries inside an added file never count.

The transcript shows a run in which \file{dangling.py} disconnects the output
\code{irq_o} of \code{u_periph}. The signal \code{irq} is read by the top module
and is now driven by nobody: \emph{left dangling}. In chapter~\ref{ch:operations},
\file{ports.py} disconnects \code{count1_o}, which goes to an output of the top
module: that is allowed.

The sign-off does not see everything. Chapter~\ref{ch:limits} lists what it misses,
for example a signal that loses the driver of only some of its bits.

\section{Post-conditions}

Each applied operation registers a \term{check} that is evaluated on the final
design. The check makes sure that a later edit has not undone the operation. A failed
check starts with the name of the operation and the target. Examples:

\begin{itemize}
\item \code{failed check: add_wire soc_top tmp: declaration not found in the final
  design}
\item \code{failed check: connect soc_top.u_rst_sync test_i: expected 'x', found 'y'}
\item \code{failed check: add_file NAME: none of its modules is instantiated}
\item \code{failed check: propagate_param WIDTH PATH: the value in use changed
  from 8 to 4}
\end{itemize}

If you see one, look at the later operations of your script (or of a later script):
one of them has changed what this operation wrote. Two operations that cancel each
other, such as \apiname{Module.add_wire} and \apiname{Module.remove_wire} on the same name,
fail in this way.

Skipped operations register no check.

\section{Expectations}

The sign-off compares the structure of the baseline and the final design, instance by
instance. A \term{change} has a kind, the path of the instance, a name and the text
before and after.

\begin{description}
\item[instance-added, instance-removed] An instance appears or goes. An added or
  removed instance shows only this change, and the same change for each instance
  below it.
\item[module-changed] The instance uses another module.
\item[param-changed] A parameter has another value at the instance.
\item[port-added, port-removed, port-changed] The module of the instance gains, loses or
  changes a port. A port is described by its direction and width, for example
  \code{input 1}.
\item[port-type-changed] The direction and width stay, but the declared type differs.
\item[connection-changed] The connection of a port has other text.
\end{description}

Every change must be covered by an expectation or a permission, and every
expectation must match a change. Operations register their own. For example,
\apiname{Module.add_port} requires \code{port-added} at every path of the module, and
permits \code{connection-changed} for the port. You add yours with \apiname{Vet.expect}
and \apiname{Vet.permit}.

\section{Name the top module}

hdl-vet needs to know the top module of your design. If you give none with
\opt{--top}, it takes the modules that nobody instantiates. There is a trap. If your
script replaces or removes the \emph{last} instance of a module, that module is no
longer instantiated, so it becomes a top-level module of the edited design. The
sign-off then reports \code{instance-added}: a new top appeared. This is not a
problem of your script. Name the top module, and the edited design has the same top
as the original. The transcript shows both. Put \code{top} in the project file so
that nobody forgets.

\lstinputlisting[style=console]{examples/console/top.txt}

\chapter{The command line}
\label{ch:cli}

The program is \cmd{vet}. It has two commands, \cmd{vet apply} and \cmd{vet init}, and
one option for itself, \opt{--version}. \cmd{python -m hdl_vet} is the same program.
The options \code{-h} and \code{--help} of \cmd{vet}, \cmd{vet apply} and \cmd{vet init}
print a short usage text. This chapter describes every option in full.

Chapter~\ref{ch:first} explains the file list, the base and the input tree.

\section{vet apply}

\cmd{vet apply} runs scripts on a design. These are its options.

\begin{description}
\item[\opt{-f}, \opt{--filelist} FILE] The top-level file list. It is required,
  unless you give \opt{--resume} or the project file names one. It replaces the
  file list and the base of the project.
\item[\opt{-o}, \opt{--outdir} DIR] The folder for the output tree. DIR must be a
  new folder or an empty folder, unless you give \opt{--force}. It must not lie inside
  the input tree, and it must not contain the input tree. It must not contain a script
  that runs or the project file.
\item[\opt{-s}, \opt{--script} FILE] A script. You can give it many times. The scripts run
  in the order given, after the scripts of the project file. A file that you name twice is
  an error.
\item[\opt{--top} MODULE] The top module. It replaces the top module of the project.
  Always name it (chapter~\ref{ch:signoff}).
\item[\opt{--base} DIR] The base of the file list (chapter~\ref{ch:first}). It replaces
  the base of the project.
\item[\opt{--dry-run}] Print the diff of the changed files and write no output tree.
  The sign-off runs. An archive that you name with \opt{--save} is still written.
\item[\opt{--force}] Replace an output folder that exists (below).
\item[\opt{--save} FILE] After the scripts end, save the session to a zip archive
  (chapter~\ref{ch:session}).
\item[\opt{--resume} FILE] Continue from a saved session. hdl-vet does not read a
  file list. Do not give \opt{-f}, \opt{--top} or \opt{--base} with it.
\item[\opt{--project} FILE] Use this project file instead of \file{hdl-vet.toml} in
  the current folder (chapter~\ref{ch:project}).
\item[\opt{--no-project}] Ignore the project file. You cannot give it together with
  \opt{--project}.
\end{description}

You must give at least one of \opt{-o}, \opt{--save} and \opt{--dry-run}. The exception
is a project file that names an output folder (chapter~\ref{ch:project}). All file names
and folder names on the command line start from the current folder.

\subsection*{What it prints}

\begin{itemize}
\item With an output tree, hdl-vet prints these lines, in this order:
  \begin{itemize}
  \item the line \code{wrote DIR}
  \item the line \code{saved FILE}, if you saved
  \item the line \code{N applied, M skipped}
  \item one line \code{changed PATH} for each changed file, with the name relative to the
    tree
  \item one line \code{added NAME} for each file that a script added
  \item the \code{note:} lines
  \end{itemize}
\item With \opt{--dry-run}, hdl-vet prints the diff on the standard output. It prints
  everything else on the error output.
\item If hdl-vet uses a project file, the first line, on the error output, is
  \code{project: PATH}.
\item If the run fails, hdl-vet prints the message on the error output.
\end{itemize}

\subsection*{Exit statuses}

This table is the one place that defines the exit statuses.

\begin{center}
\begin{tabular}{@{}lp{0.75\textwidth}@{}}
\toprule
Status & Meaning \\
\midrule
0 & The run worked. Without \opt{--dry-run}, hdl-vet writes the output tree. With \opt{--save}
  and no output, the run ends with 0 also when the sign-off failed so far
  (chapter~\ref{ch:session}). \\
1 & The sign-off failed (chapter~\ref{ch:signoff}). The report is on the error output, and
  hdl-vet writes no output tree. This holds with \opt{--dry-run} too. \\
2 & hdl-vet cannot do the request: a wrong command line, a file that it cannot read,
  a design that it cannot load at the start, a script that fails, a refusal of an
  operation, or drift. It writes nothing. \\
\bottomrule
\end{tabular}
\end{center}

Messages of the command line have two forms. The first form is for a command line that is
wrong in a way that the argument parser sees. Examples are an unknown option, a missing
value, and no command at all. In this case \cmd{vet} prints a usage line and then
\code{vet: error: ...}, and exits with status 2. For every other error, \cmd{vet} prints
one line, \code{vet: ...}, and exits with status 2.

\lstinputlisting[style=console]{examples/console/cli.txt}

In the transcript, look at these things:

\begin{itemize}
\item \cmd{vet --version} prints the program name and the version (here \code{<version>}).
\item The file list \file{flow.f} is outside the folder of the sources, so the
  command gives \opt{--base} \code{design}. The output tree has the file list at its top.
\item hdl-vet refuses the second run, because \file{out/flow} exists. The message names
  \opt{--force}. The third run uses it.
\item The transcript has five errors, each with exit status 2:
  \begin{itemize}
  \item no file list
  \item no output
  \item a file list that is not there
  \item a top module that is not in the design
  \item a script that cannot be read
  \end{itemize}
\end{itemize}

\subsection*{What replacing an output folder means}

\opt{--force} replaces the output folder. hdl-vet builds the new tree next to the
output folder. Then it renames the new tree to the name of the output folder. So a failed
run never leaves half a tree. This is the one place that describes the limits.

\begin{itemize}
\item With \opt{-o}, \opt{--force} replaces \emph{any} folder that you name,
  and this includes a folder that holds files of yours. Check the name.
\item It never replaces a folder inside the input tree, the input tree itself, or a
  folder that holds the input tree. It never replaces a folder that holds a script
  of the run or the project file.
\item If the output folder comes from the project file, \opt{--force} replaces it in
  two cases only:
  \begin{itemize}
  \item The folder is empty.
  \item The folder is the output of an earlier run and holds nothing else.
  \end{itemize}
  In the second case, \file{hdl-vet.json} must be a manifest of hdl-vet with a list of
  files. Every file in the folder must be that manifest or be in the list. Otherwise
  hdl-vet refuses. The message names a file that is not in the list.
  Output trees from before the list existed have no such list, so \opt{--force} refuses
  them.
\end{itemize}

\section{vet init}

\cmd{vet init DIR -f FILELIST} writes a starter project into \file{DIR}
(chapter~\ref{ch:first} and chapter~\ref{ch:project}).

\begin{description}
\item[\opt{-f}, \opt{--filelist} FILE] The file list of the design. It is required.
\item[\opt{--top} MODULE] The top module. If you give none, \cmd{vet init} takes the only
  top-level module of the design. If the design has several, \cmd{vet init} lists them
  (at most ten, the biggest hierarchy first) and asks for \opt{--top}.
\item[\opt{--base} DIR] The base of the file list.
\end{description}

\cmd{vet init} loads the design first, so it refuses a design that has errors. It
refuses a \file{DIR} that exists and is not empty. It refuses a \file{DIR} inside the
input tree, because the next run of the generator can delete it. When it refuses, it writes
nothing. It prints \code{created DIR}, exactly as you typed it, and one line
\code{wrote PATH} for each file.

\chapter{The project file}
\label{ch:project}

A project file names the design, the scripts and the output folder, so that you can
run \cmd{vet apply} with no arguments. The file is called \file{hdl-vet.toml}. It is
a TOML file. \cmd{vet init} writes one for you (chapter~\ref{ch:first}).

\section{The keys}

\begin{description}
\item[\code{[design]} \code{filelist}] The top-level file list. It is required when
  the project gives the design.
\item[\code{[design]} \code{top}] The top module. It is optional, but name it
  (chapter~\ref{ch:signoff}).
\item[\code{[design]} \code{base}] The base of the file list. It works as \opt{--base}.
  You can give it only together with \code{filelist}.
\item[\code{[apply]} \code{scripts}] A list of script files, in the order in which they
  run. It is optional.
\item[\code{[apply]} \code{output}] The output folder. It is optional.
\end{description}

File names in the project file start from the folder of the project file, and use
\code{/}. hdl-vet takes an absolute name as it is. Here is the file that \cmd{vet init}
wrote in chapter~\ref{ch:first}:

\begin{itemize}
\item \code{filelist = "../design/design.f"}
\item \code{top = "soc_top"}
\item \code{scripts = ["scripts/01_example.py"]}
\item \code{output = "out"}
\end{itemize}

hdl-vet reads the file strictly. These things are an error:

\begin{itemize}
\item an unknown table or key
\item a value of the wrong type
\item an empty string
\item a nested table
\item text that is not TOML
\end{itemize}

The error names the file and the key. hdl-vet guesses nothing. The file must be UTF-8.

\section{Where it is looked for}

\cmd{vet apply} uses \file{hdl-vet.toml} in the current folder, or the file that
\opt{--project} names. It does not look in parent folders. \opt{--no-project} ignores
the file. These things are an error:

\begin{itemize}
\item \opt{--project} and \opt{--no-project} together
\item an \opt{--project} file that does not exist
\item an \opt{--project} name that is a folder
\end{itemize}

Whenever hdl-vet uses a project file, the first line on the error output is
\code{project: FILE}, with the full name. A project never runs without notice.

\section{What the command line replaces}

\begin{itemize}
\item \opt{--top} replaces \code{top} of the project.
\item \opt{-f} replaces \code{filelist} and \code{base} of the project together,
  because a base belongs to its file list. \opt{--base} replaces only the base.
\item \opt{-s} adds scripts. The scripts of the project run first, then the scripts of
  the command line. When you give \opt{-f} and the project has scripts, hdl-vet prints
  this note: \code{note: the project's scripts run first: ...}.
\item hdl-vet uses \code{output} of the project only when you give none of \opt{-o},
  \opt{--save} and \opt{--dry-run}. \cmd{vet apply --save FILE} saves and writes no
  tree, as without a project.
\item With \opt{--resume}, hdl-vet does not use the design of the project. The saved
  session has its own design. hdl-vet compares the scripts of the project with the
  scripts that the session applied (chapter~\ref{ch:session}).
\end{itemize}

Everything else is as without a project. All checks of \opt{--save} and
\opt{--resume} apply to the combined settings. The project file counts as a file that
\opt{--save} must not replace. Paths on the command line start from the current
folder, even when \opt{--project} names a file elsewhere. The \code{wrote} line shows
an absolute name when the output folder comes from the project file. If you give the
folder with \opt{-o}, the line shows it as you typed it.

\section{Safety}

When hdl-vet writes an output tree, it replaces the output folder. So the output folder,
wherever it comes from, must not contain a script of the run or the project file. The
\code{output} of a project must not be the folder of the project. It must not contain
this folder. Chapter~\ref{ch:cli} gives the rules for \opt{--force} on an output folder
from the project file.

\cmd{vet init} also writes these things:

\begin{itemize}
\item a \file{.gitignore} with \code{/out/} and \code{*.session.zip}
\item a \file{README.md} with the commands
\item a \file{cells} folder with a note about where to keep files for
  \apiname{Vet.add_file}
\end{itemize}

\lstinputlisting[style=console]{examples/console/project.txt}

The transcript shows these steps:

\begin{itemize}
\item \cmd{vet init} runs twice. The second time, the folder is not empty, and hdl-vet
  refuses.
\item A run in the project with an extra script. The script \file{skipped.py} runs
  after \file{01_example.py}, and the header line lists both.
\item \cmd{vet apply} writes \file{out}. hdl-vet refuses the same command again, because
  \file{out} exists. With \opt{--force}, hdl-vet replaces the folder. This works because
  the old \file{out} is an earlier output of hdl-vet with nothing else in it.
\item \opt{-f} on the command line. hdl-vet prints the note about the scripts of the
  project.
\item \opt{--project} from another folder, then the options that conflict, then a run
  without a project and without a file list.
\end{itemize}

\chapter{Saving and resuming a session}
\label{ch:session}

On a big design, loading takes time, and you may want to try a new script again and
again after the same first two. A \term{session} is the state of a run: the text of every
file with the edits so far, the expectations and checks of the operations, the log, the
notes, the added files and the scripts that ran. You can save it and continue later,
without the input tree and without running the first scripts again.

\section{Saving}

\opt{--save} \code{FILE} writes the session to a zip archive after all scripts have run. The
archive holds plain data, and nothing in it is code: loading it never runs anything.
You can open it with any zip tool. It has \file{session.json} (the index),
\file{original/} (every source file as it was loaded), \file{current/} (the files that
differ) and \file{added/} (the files that scripts added). The same state gives the same
bytes. An existing archive is replaced.

\begin{itemize}
\item Without \opt{-o}, hdl-vet saves, prints \code{saved FILE} and the summary, and then
  one line, \code{sign-off so far: ok} or \code{sign-off so far: FAILED}. After
  \code{FAILED} the report follows on the error output. The exit status is 0 in both
  cases: saving a state that does not pass is not an error. The state is saved as it is.
\item With \opt{-o} or \opt{--dry-run} too, the archive is written once every script has
  run, before the sign-off. The run then goes on as usual: it can still exit with
  status 1, and the archive exists either way. The line \code{sign-off so far} is not
  printed.
\item A script that fails writes no archive.
\item hdl-vet refuses a \opt{--save} file that lies inside the input tree or inside the
  \opt{-o} folder, or that is a script, the project file or a file list that the design
  reads: the top-level one, or one that it reads with \code{-f} or \code{-F}, wherever
  that list lies. A resumed session refuses the same file lists.
  The scripts are then not run: \code{apply} is not called. The code at the top level of
  every script has already run by then, because hdl-vet reads each script before it
  checks \opt{--save}.
\end{itemize}

\section{Resuming}

\opt{--resume} \code{FILE} continues from a saved session instead of reading a file list. Do
not give \opt{-f}, \opt{--top} or \opt{--base}: the session has them, and hdl-vet refuses
them. You add scripts with \opt{-s}, and you may save again with \opt{--save}; that
replaces the old archive when it names the same file.

\begin{itemize}
\item A script that the session has applied already (same name, same content) is
  refused.
\item The summary counts only the operations since the resume, and only the notes since
  the resume are printed.
\item The sign-off covers every script, the old ones and the new. The header line and the
  manifest list all scripts, the saved ones first. If you save the state and resume it
  with one more script, the output tree is byte for byte the tree of one run of all the
  scripts.
\item The session reports and lays out files exactly as the first run did. It needs the
  same folders: it records where the input tree was.
\end{itemize}

With a project file, the project's scripts are compared with the saved list from the
start, by name and content. If the project's script number $n$ is not the saved script
number $n$, hdl-vet stops and says that the project has changed since the save.
\opt{--no-project} resumes regardless. Project scripts that the session does not have
yet run first, then the scripts of \opt{-s}.

\lstinputlisting[style=console]{examples/console/session.txt}

\section{What is and is not detected}

A saved session is used as it is. hdl-vet does not check whether the file list, the
sources or the scripts changed after the save. In the transcript, the command \cmd{cp}
puts a regenerated \file{soc_top.sv} into the design after the save. The resume does not
notice. It continues from the text that it saved, so the output has no \code{busy_o}.
If your generator ran again, start from the beginning.

What hdl-vet does detect:

\begin{itemize}
\item A script that has been applied already (above).
\item An archive that is not a session archive, is damaged, or has another format
  version. The message names the file.
\item A project whose scripts differ from the saved ones (above).
\item \opt{-f}, \opt{--top} or \opt{--base} together with \opt{--resume}.
\end{itemize}

A scope or a group that is open cannot be saved (chapter~\ref{ch:scripts}). Nothing about
scopes is kept in a session.

\chapter{Using hdl-vet from Python}
\label{ch:python}

\cmd{vet apply} is a short program on top of a Python interface. You can use the
interface in your own code, for example in a test. This chapter is a reference. It has
no example script, because the command line runs the scripts of the other chapters in
the same way.

\section{Building a \texttt{Vet}}

The command line does this:

\begin{enumerate}
\item It builds a session with \code{hdl_vet.session.Session(filelist, top=None,
  editable=None, base=None)}. The session loads the baseline design. It raises
  \code{SessionError} when the baseline has errors, or when a source file already
  carries the hdl-vet header. It raises \code{FilelistError} for a file list that cannot be
  read, and \code{AdapterError} when the compiler cannot load the design.
\item It wraps the session in \code{hdl_vet.ops.Vet(session)}.
\item It calls the function \code{apply(vet)} of each script.
\item It calls \apiname{Vet.finish} with the names and hashes of the scripts. If the
  result is \code{ok}, it writes the tree with \code{hdl_vet.output.write_tree}.
\end{enumerate}

The errors of the operations are \code{hdl_vet.ops.VetError}, and its subclass
\code{DriftError}. A script may catch them (chapter~\ref{ch:operations} has an
example). \code{vet.log} is the list of \code{OpRecord(op, target, status)} with the
status \code{"applied"} or \code{"skipped"}, and \code{vet.notes} is the list of notes.

\section{The calls}

\begin{entry}{\sig{Vet.check}}
\item[Returns] The \code{Result} that \apiname{Vet.finish} would return now, and leaves the
  session as it was: text, queue and operations stay, and you can go on. Pass the
  \code{scripts} (pairs of name and hash) that should appear in the header.
\item[Refused] Inside a group or a scope, and after \apiname{Vet.finish}.
\end{entry}

\begin{entry}{\sig{Vet.finish}}
\item[Returns] The \code{Result} of the sign-off. It applies the queue, adds the header
  lines, loads the edited design, compares, and evaluates the checks. A session can be
  finished once. The header lists the scripts that the session applied before, then those
  that you pass.
\item[Result] \code{ok}; \code{new_diagnostics}; \code{dangling}; \code{changes};
  \code{unexpected_changes}; \code{unmatched_expects}; \code{failed_checks};
  \code{new_files}; \code{load_error}; \code{files} (new bytes of each changed file);
  \code{added} (the files that scripts added); \code{diff} (the unified diff); and the
  method \code{report()}, which prints the lines of chapter~\ref{ch:signoff}.
\end{entry}

\begin{entry}{\sig{Vet.save}}
\item[Writes] The session to the archive \code{path}. It first applies and loads what is
  queued. \code{scripts} are the pairs of name and hash of the scripts applied since the
  session began or resumed. The session stays usable.
\item[Refused] Inside a group or a scope, after \apiname{Vet.finish}, when the queued edits
  cannot be loaded, and when \code{path} is inside the input tree or is a file that the
  session owns (the file list, a source of an added file).
\end{entry}

\begin{entry}{\sig{Vet.restore}}
\item[Returns] A new \code{Vet} in the state saved in \code{path}. It is a class
  method: \code{Vet.restore(path)}. It raises \code{SessionError}, with the file named,
  when the file is not a session archive, is damaged, has another format version, or
  names a file outside its layout root.
\end{entry}

\begin{entry}{\sig{Vet.text}}
\item[Returns] The current text of a file of the design, as bytes (chapter~\ref{ch:operations}).
\end{entry}

\begin{entry}{\sig{Vet.design}}
\item[Returns] The loaded design with every queued edit applied. It loads the design
  again if something is queued. Its type is the adapter's \code{Design}, which wraps
  the compiler. Only \code{hdl_vet.slang_adapter} imports \code{pyslang}.
\end{entry}

\begin{entry}{\sig{Vet.module_paths}}
\item[Returns] The sorted hierarchical paths of the module in the baseline design, as a
  list; empty for a module that is not instantiated.
\end{entry}

\begin{entry}{\sig{Vet.open_block}}
\item[Returns] \code{"scope 'NAME'"} when a scope is open, \code{"group"} when a group
  is open, else \code{None}. The command line uses it to refuse a script that leaves a
  block open.
\end{entry}

\chapter{Limits and known problems}
\label{ch:limits}

This chapter lists what hdl-vet does not do, and what is known to go wrong. Every
item was checked against the code of this version. Items with a transcript were run.
Items that say so were not checked by a test.

\section{What the sign-off does not see}

\begin{itemize}
\item \textbf{Logic.} The sign-off compares the structure (instances, ports,
  parameters, connections) and the diagnostics of the compiler. A raw edit that changes
  only logic passes without an \code{expect}. The first transcript below changes
  \code{count_o + 1'b1} into \code{count_o + 2'd2} in \code{timer}, and the sign-off
  says nothing.

  Use the operations where you can. They never write such an edit. If you
  use \apiname{Vet.splice}, you are responsible for the logic. The same holds for
  declarations, because \apiname{Module.add_wire} has no structural expectation.
\item \textbf{Partly undriven signals.} A signal can lose the driver of only some of
  its bits. This can happen through an operation that is not a checked operation (here
  \apiname{Instance.disconnect}), or through a splice. The sign-off passes such a run.

  The second transcript disconnects the interrupt of the second timer. Bit 1 of
  \code{irq_o} of \code{periph_wrap} has no driver, and the top module reads it. The
  sign-off accepts the run. Checked operations see this case (chapter~\ref{ch:operations},
  \file{remove\_timer.py}).
\item \textbf{Operations that cancel each other.} If one operation undoes another
  operation of the same script, the sign-off fails. An example is
  \apiname{Module.add_wire} and then \apiname{Module.remove_wire} of the same name. The
  check of the first operation does not find its result. The third transcript shows it.
  Remove both operations from the script.
\item \textbf{Gate primitives, user-defined primitives and checkers.} The sign-off does
  not compare them. A splice that changes \code{and g1 (z, a, b)} into
  \code{and g1 (z, a, a)} passes. hdl-vet writes the output tree, with exit status 0.
  The gate primitive was run. The other two follow from the same rule and were not checked
  by a test.
\item \textbf{More that the dangling check does not see.} The dangling check does not see
  these things:
  \begin{itemize}
  \item a parameter change that switches off a driver in a generate branch
  \item an inout port
  \item a new reader of a signal that had no driver already in the baseline
  \end{itemize}
  These cases were not checked by a test. The specification of the check lists them.
\item \textbf{One configuration.} The sign-off sees the design with the defines of the
  file list. An operation refuses a list that holds a conditional compilation directive,
  because it cannot check another configuration.
\end{itemize}

\lstinputlisting[style=script]{examples/scripts/logic_change.py}
\lstinputlisting[style=script]{examples/scripts/partial.py}
\lstinputlisting[style=script]{examples/scripts/cancel.py}

\lstinputlisting[style=console]{examples/console/limits.txt}

\section{What the operations refuse}

hdl-vet refuses to edit these things, and it prints a message. You can edit them by hand,
or with a raw splice.

\begin{itemize}
\item instances in generate blocks and instance arrays, and instances that come from a
  macro
\item hierarchical paths with an index
\item connection lists with positional or \code{.*} connections, and parameter lists
  with positional overrides
\item modules with a non-ANSI port list (for port operations), and module headers with
  a macro
\item ports and parameters inside conditional compilation directives
\item modules that are not instantiated under the top module
\item modules that come from an added file. hdl-vet never edits an added file.
\item interface ports. \apiname{Vet.punch_port} refuses them and
  \apiname{Instance.replace_module} refuses a module that has one. There is no operation
  that adds an interface port.
\item a replacement (\apiname{Instance.replace_module}) by a module whose port has another
  direction, width or type than the port that it replaces. There is no option to accept
  such a difference.
\item type parameters. \apiname{Instance.replace_module} does not carry an override of a
  type parameter to the new module. This holds even if the new module has a type parameter
  of that name. The message says that the new module has \emph{no overridable parameter}.
  Edit such an instance by hand.
\end{itemize}

\section{Known problems}

\begin{itemize}
\item \textbf{Escaped names.} \apiname{Module.remove_instance} cannot remove an instance
  with an escaped identifier, such as \texttt{\textbackslash u/esc}. If you give the name
  with the backslash, hdl-vet skips the call without a message, although the instance
  exists. If you give the name without the backslash, hdl-vet applies the call, but the
  sign-off fails with
  \code{unexpected change: instance-removed}.
\item \textbf{Names that the operations do not check.} \apiname{Module.add_port} and
  \apiname{Module.add_param} do not check the name that you give. They write
  \code{"a b"} or a keyword into the text, and the sign-off fails with errors of the
  compiler or with \code{the edited design cannot be loaded}. \apiname{Module.add_wire}
  does check its name.
\item \textbf{Time.} An operation that reads a list that an earlier operation edited
  loads the whole design again. On a large design, a script with many such operations takes
  a long time. The load count rose after a \code{connect} that follows an
  \code{add_port}. The times were not measured.
\item \textbf{Paths after a replacement.} After \apiname{Instance.replace_module}, a
  hierarchical operation (\apiname{Vet.punch_port}, \apiname{Vet.propagate_param}) still
  resolves its path against the original design. It edits the old module, and the sign-off
  reports the result as errors or failed checks. Do such operations before the
  replacement, or in another run.
\item \textbf{Scopes and shared modules.} hdl-vet refuses \apiname{Vet.punch_port} and
  \apiname{Vet.propagate_param} in a scope whose module the design instantiates more than
  once. The scope takes the name of a module, not an instance path. The check for
  references that reach outside a scope does not see a name in code that the compiler does
  not elaborate. Its message names only one reference when a block adds several.
\item \textbf{Added files.} A relative file name in \apiname{Vet.add_file} starts from the
  current folder, not from the script. Use \code{Path(__file__).resolve().parent} as in the
  example scripts.
\item \textbf{Saved sessions of an older version.} A session that you saved before
  hdl-vet recorded the nested file lists protects only the top-level file list when you
  resume it. With such a session, \opt{--save} onto a nested file list outside the input
  tree overwrites that list.
\item \textbf{Replacing a folder.} \opt{--force} with \opt{-o} replaces any folder that is
  not part of the input tree. This includes a folder with your own files in it. The rules
  are in chapter~\ref{ch:cli}.
\item \textbf{Old output trees.} An output tree that an older version wrote has no list
  of files in \file{hdl-vet.json}. \opt{--force} refuses it when the folder comes from the
  project file.
\end{itemize}


\end{document}
